\documentclass[AMA,Times1COL]{WileyNJDv5}

\articletype{Research Article}%

\received{Date Month Year}
\revised{Date Month Year}
\accepted{Date Month Year}
\journal{Statist.\ Med.}
\volume{00}
\copyyear{2026}
\startpage{1}

\raggedbottom

\usepackage{mathtools}
\usepackage{amssymb}
\input{numbers_macros}  % GENERATED from numbers_registry.csv by scripts/make_numbers_macros.R — every \numval{} is a registry lookup

\graphicspath{{../figures/}}

\begin{document}

\title{Quantifying Effect Modifier Similarity for Regional Pooling in Multi-Regional Clinical Trials}

\author[1]{Author One}

\author[2]{Author Two}

\author[1]{Author Three}

\authormark{AUTHOR ONE \textsc{et al.}}
\titlemark{QUANTIFYING EFFECT MODIFIER SIMILARITY FOR REGIONAL POOLING}

\address[1]{\orgdiv{Department Name}, \orgname{Institution Name}, \orgaddress{\state{State}, \country{Country}}}

\address[2]{\orgdiv{Department Name}, \orgname{Institution Name}, \orgaddress{\state{State}, \country{Country}}}

\corres{Corresponding Author, Department Name, Institution Name, Address. \email{corresponding@institution.edu}}

% Abstract rewritten 2026-09-27 (Terminal pass, decision round under Tak's instruction "判断が要る点を議論して判断して"): <=250 words, no citations (Statistics in Medicine); every claim traced to results/verify_study2.log and results/verify_study1_app.log.
% Abstract revised 2026-10-12 (D-20261012-2, Louis M1/M2/m2; individual veto open until submission): the MCSE count of D-20261011-11 judgement 1(a) carried into the abstract; "three by construction" attached to 11 of 13; "about 0.4 SD or more" (S5 is 0.399 SD); "equal regional sizes" restored.
\abstract[Abstract]{%
% Cut 2026-10-15 (D-20261015-1, needs Tak's veto): abstract revised last, after the body cuts: 'regional differences in the average treatment effect are at most Delta_max' -> 'regional average treatment effects differ by at most Delta_max ... on a risk- or mean-difference scale' (restores the collapsible-scale qualifier, guard LR3-G6); every other phrase checked against the shortened text (selection simulation with matches supplied, k of ten cells, covExtE6-covExtE5 at the application's operating point, null floor) and unchanged; the coverage sentence carrying four guards untouched
The ICH E17 guideline describes regional pooling in multi-regional clinical trials (MRCTs) in terms of the similarity of effect modifier distributions but names no measure, and to our knowledge no distance between regional effect modifier distributions has been evaluated for it. We quantify how two regions' distributions of each candidate effect modifier differ by the Wasserstein-1 ($W_1$) distance in original units. If the conditional average treatment effect is common to all regions, depends only on that modifier and is Lipschitz with clinical slope $L_{\text{clinical}}$, regional average treatment effects differ by at most $\Delta_{\max} = L_{\text{clinical}} W_1$ on a risk- or mean-difference scale; when $L_{\text{clinical}}$ is unknown, the smallest slope $L^*$ at which the distance could produce a clinical margin is compared with a sponsor's bound. In a selection simulation (matches supplied), the standardized mean difference stayed at chance up to 2,000 patients per region in some configurations; $W_1$ and the Kolmogorov--Smirnov (KS) statistic reached the target (AUC 0.9) in every configuration. For true distances of about \numval{minW1SdS3toS7} to \numval{maxW1SdS3toS7} SD (normal reference; equal $n$), bootstrap coverage from 100 patients per region was \numval{covRangeS3toS7N100up2dpAbs}, below 0.95 beyond Monte Carlo error in \numval{nCellsCovBelow2MCSE} of ten cells (\numval{covExtE6}--\numval{covExtE5} at the application's operating point). In a GUSTO-I-based example, partner rankings on age and systolic blood pressure differed; a null floor marked distances not distinguishable from no difference. The framework quantifies differences; the pooling judgment remains clinical.}

\keywords{multi-regional clinical trial; ICH E17; effect modifier; Wasserstein distance; regional pooling; distributional similarity}

\jnlcitation{\cname{%
\author{Author One},
\author{Author Two}, and
\author{Author Three}}.
\ctitle{Quantifying effect modifier similarity for regional pooling in multi-regional clinical trials.} \cjournal{Statist.\ Med.} \cvol{2026;00:1--23}.}

\maketitle

\renewcommand\thefootnote{}
\footnotetext{\textbf{Abbreviations:} AMI, acute myocardial infarction; AUC, area under the receiver operating characteristic curve; CATE, conditional average treatment effect; CDF, cumulative distribution function; CI, confidence interval; CRP, C-reactive protein; CV, coefficient of variation; EDF, empirical distribution function; GUSTO, Global Utilization of Streptokinase and t-PA for Occluded coronary arteries; ICH, International Council for Harmonisation; IPD, individual patient data; KL, Kullback--Leibler; KS, Kolmogorov--Smirnov; MRCT, multi-regional clinical trial; MRI, magnetic resonance imaging; NMPA, National Medical Products Administration; RMSE, root mean squared error; RWE, real-world evidence; SBP, systolic blood pressure; SD, standard deviation; SMD, standardized mean difference.}

\renewcommand\thefootnote{\fnsymbol{footnote}}
\setcounter{footnote}{1}

%==============================================================================
\section{Introduction}\label{sec:intro}
%==============================================================================

% Cut 2026-10-15 (D-20261015-1, needs Tak's veto): clause 'and methods for assessing the consistency of treatment effects across their regions have developed with them' -> dropped (pure deletion of the clause)
Multi-regional clinical trials (MRCTs), conducted across multiple countries or regulatory regions under a single protocol, have become a common design in global pharmaceutical development.\cite{chen2010,quan2010} The International Council for Harmonisation (ICH) E17 guideline, adopted in 2017, established principles for planning and designing MRCTs to evaluate whether treatment effects can be regarded as generalizable across the target population.\cite{iche17}

% Cut 2026-10-14 (D-20261014-1 I01, F): CLAIMS C-01, C-18
% Cut 2026-10-15 (D-20261015-1, needs Tak's veto): PREVENT paragraph (about 150 words: binary characteristics, attribution of the Japanese inconsistency to assignment imbalance, 'the step a distance would quantify') -> one factual sentence appended to this paragraph (insurance A; the fact is the one the removed paragraph stated, source matsushima2024)
Regulatory authorities in Japan and China examine the consistency of treatment effects in regional subpopulations,\cite{matsushima2024,song2025} yet individual regional sample sizes in global MRCTs are often insufficient for that purpose. Regional pooling, in which regions with similar patient characteristics are grouped for analysis, is a key strategy for enhancing the assessment of regional consistency, and the ICH E17 guideline describes pooling decisions in terms of the similarity of effect modifier distributions.\cite{iche17} In the PREVENT trial of secukinumab, a workshop report describes countries with a similar distribution of effect modifiers being clustered post hoc into five pooled regions on binary characteristics (CRP and MRI positivity, sex).\cite{matsushima2024}

% Cut 2026-10-14 (D-20261014-1 I02, F): CLAIMS C-19
% Cut 2026-10-15 (D-20261015-1, needs Tak's veto): the worked example ('a region with predominantly younger patients would show larger benefits than one with predominantly older patients because the patient mix differs, not the drug') and 'such as age and disease severity scores' after 'continuous effect modifiers' -> dropped; 'Where it exists' -> 'therefore'
An effect modifier is a baseline patient characteristic, such as age, disease severity or a genetic marker, for which the treatment benefit differs across subgroups. Regions with different patient compositions may therefore observe different average treatment effects even if the drug works identically at the individual level, so the similarity of the full effect modifier distribution, not only its average level, matters for regional pooling. We focus on continuous effect modifiers; categorical effect modifiers raise distinct considerations and are outside the present scope.

% Cut 2026-10-14 (D-20261014-1 I03, F): CLAIMS C-02 (context only)
% Cut 2026-10-14 (D-20261014-1 H1, F): CLAIMS C-02 (context only); Harvey menu 1

% Edited 2026-09-27 (decision round under Tak's instruction): Song sentence corrected to what Song et al. 2025 state (pooling strategy an area of remaining ambiguity; cluster analysis described, p.4); next paragraph's "This absence of quantitative tools" replaced, since this paragraph itself describes quantitative clustering proposals.
% Cut 2026-10-14 (D-20261014-1 I04, F): CLAIMS C-20
% Cut 2026-10-14 (D-20261014-1 I05, F): CLAIMS C-02
% Cut 2026-10-14 (D-20261014-1 I06, F): CLAIMS C-03
% Cut 2026-10-15 (D-20261015-1, needs Tak's veto): Komiyama chapter description (representative value per region, standardization, Lasso-weighted Euclidean distance, 'selection of the distance measure is one of the key decisions') -> 'introduces, as one of many possible approaches, cluster analysis of regions on region-level summaries ... and does not evaluate it'; 'when several effect modifiers are considered' (Song) and the binary-characteristic clause of the first gap dropped
Despite the regulatory importance of effect modifier distributional similarity, current practice lacks a standardized, clinically calibrated methodology. The ICH E17 guideline provides no specific metric, threshold, or statistical procedure for determining when distributions are ``similar enough.'' Song et al., on the NMPA's implementation of ICH E17, identify the pooling strategy as an area where ambiguity remains and describe clustering populations on distances defined from the correlation between effect modifiers and endpoints.\cite{song2025} The pooling-strategy chapter of a reference volume on multi-regional trials under ICH E17 introduces, as one of many possible approaches, cluster analysis of regions on region-level summaries of the candidate effect modifiers, and does not evaluate it.\cite{komiyama2024} Two gaps remain. First, to our knowledge no distance between regional effect modifier distributions has been proposed and evaluated for this purpose: a single summary per modifier leaves two regions that agree on it but differ in spread, shape, or tail behaviour at distance zero, and any finite set of summaries requires deciding in advance which distributional features are relevant. Second, the chapter calls for a judgment of whether differences in the impact of effect modifiers across pooled regions are ``clinically meaningful or not'' but supplies no metric to operationalize that judgment.

% Paragraph added 2026-10-12 (D-20261012-3, EXT M1 / REF M1; individual veto open until submission): prior literature positioned (sources grounded in review_response/LIT_POSITIONING_20261012.md); novelty restated as formulation, reverse calibration and evaluation, not the inequality
% Cut 2026-10-15 (D-20261015-1, needs Tak's veto): the generalizability-literature clause (dahabreh2020) dropped; 'the same bound on a difference of expectations under a Lipschitz function' and the propensity-score / generalizability-index / trimmed-Mallows sentences shortened; shen2018, redko2017 attributions, 'None of these converts ...' and the contribution sentence unchanged -> deleted (no replacement wording)
The ingredients are not new in themselves. The inequality on which the clinical calibration below rests is the Kantorovich--Rubinstein duality; the same bound under a Lipschitz function is used in the optimal-transport analysis of domain adaptation,\cite{shen2018} and a Wasserstein-1 bound of the same family holds under a reproducing-kernel condition.\cite{redko2017} The similarity between a trial sample and a target population has been summarized by propensity-score-based metrics\cite{stuart2011} and a generalizability index,\cite{tipton2014} and the similarity of two distributions within a trimmed Mallows-distance neighbourhood has been framed as an equivalence-testing problem.\cite{munk1998} None of these converts the difference between two regions' distributions of an individual effect modifier into a bound on the regional effect difference whose constant is a clinical slope, or evaluates such a distance for regional pooling. The contribution of this paper is the formulation for regional pooling in MRCTs, the reverse calibration $L^*$ when the slope is unknown, and the evaluation of the estimator's identification and estimation properties.

% Cut 2026-10-15 (D-20261015-1, needs Tak's veto): PREVENT paragraph deleted (see the marker above the Regulatory-authorities paragraph; one sentence kept there) -> deleted (no replacement wording)

% Paragraph reviewed and approved by Tak 2026-08-31 (Option B of the paragraph
% review): intent argument replaced by the structural statement, "cannot detect"
% aligned with "captures only" (Section 2.1 / Discussion), and the closing
% clause cut as redundant with the end of the preceding paragraph.
General-purpose comparison tools such as the standardized mean difference are routinely applied to baseline covariates,\cite{austin2011} but the standardized mean difference captures only differences in location: regions can match on means yet differ in variance or distributional shape. When effect modification is non-linear, these features directly influence regional average treatment effects.

% Paragraph reviewed and approved by Tak 2026-09-01 (Option B of the paragraph
% review): "simultaneously" replaced by the paragraph-7 wording (variance and
% shape as well as location); the closing "theoretically grounded" sentence cut
% -- the grounding (KR duality, heterogeneity bound) is Section 2's job.
% Cut 2026-10-14 (D-20261014-1 I08, R3d): CLAIMS C-04 (area-between-CDFs is stated in Sec. 2.2 and Fig. 1)
% Cut 2026-10-14 (D-20261014-1 I10, R3d): CLAIMS C-19, C-04
This paper addresses these gaps by proposing a quantitative framework for assessing effect modifier distributional similarity across regions, based on the Wasserstein-1 ($W_1$) distance between two cumulative distribution functions, retained in the original measurement units of the effect modifier. We estimate $W_1$ with bootstrap confidence intervals and link it to clinical interpretation through a calibration framework.

% Cut 2026-10-15 (D-20261015-1, needs Tak's veto): roadmap paragraph 'The remainder of this paper is organized as follows ...' deleted (pure deletion; no later sentence refers to it) -> deleted (no replacement wording)

%==============================================================================
\section{Methods}\label{sec:methods}
%==============================================================================

% Cut 2026-10-14 (D-20261014-1 M01, F): CLAIMS C-19
Let the conditional average treatment effect (CATE) be denoted $\tau(x) = E[Y(1) - Y(0) \mid X = x]$, where $X$ is the effect modifier. When this function is non-constant, the average treatment effect observed in region $r$ depends on the distribution $F_r$ of the effect modifier in that region:
\begin{equation}
\bar{\tau}_r = \int \tau(x) \, dF_r(x).
\label{eq:regional_ate}
\end{equation}
% Cut 2026-10-15 (D-20261015-1, needs Tak's veto): 'and everything built on it' and 'marginal odds and hazard ratios are not averages of their conditional counterparts and' dropped; the collapsibility requirement and its three examples kept -> deleted (no replacement wording)
Assessing pooling suitability therefore requires a measure of distributional difference between regions. Equation~(\ref{eq:regional_ate}) presupposes an effect scale on which the regional effect is the average of the conditional effects, that is, a collapsible contrast such as a risk difference, a mean difference or a difference in restricted mean survival time; marginal odds and hazard ratios fall outside Proposition~\ref{prop:heterogeneity} unless the regional effect is first defined by standardization on a collapsible scale.
% Sentence added 2026-10-12 (D-20261012-3, EXT M2 / REF Minor 8): collapsibility stated.

% Cascade edit 2026-09-27 (Tak: "supplement はあまり付けたくない"): Supplements A and B dropped (A's cancellation is stated in Section 2.4; B was internal notes, main text cites the sources for the bounds); the real-data survey then also dropped as a supplement (Tak: a check made while writing the paper, not a supplement); the two Discussion statements it supported are now carried by the GUSTO-I application (Table 5; results/app_all_methods.log, results/gusto_r8_w1_per_pair.csv). 2026-10-11 (D-20261011-11 judgement 4): the survey is cited in one Discussion sentence with generated numbers, not as a supplement.
% Rewritten 2026-09-27 (Tak): RV1-RV3 removed; Komiyama et al. introduce one example approach, not an evaluated method (chapter 4.6.1: "one of the approaches is introduced"), so it is not compared. Measures: W1, KS, SMD.
\subsection{Existing Approaches and Their Limitations}\label{sec:existing}

% Cut 2026-10-15 (D-20261015-1, needs Tak's veto): SMD display equation and KL display equation (eq:smd, eq:kl; neither is reused) -> SMD defined only by name (the 'difference in sample means divided by the pooled standard deviation' clause dropped), KL named without its formula; 'where X_r denotes ... S_pooled' dropped; Cramer-von Mises / Anderson-Darling named no more; the two paragraphs joined; 'asymmetric, infinite on non-overlap' stays in the next paragraph
The most widely used measure for comparing covariate distributions across groups is the standardized mean difference (SMD),\cite{austin2011} which by construction captures only differences in location: $N(50, 5^2)$ and $N(50, 15^2)$ yield $\text{SMD} = 0$ despite a threefold difference in spread. Distributional distance measures compare distributions beyond summary statistics: empirical distribution function (EDF) statistics such as the Kolmogorov--Smirnov (KS) statistic, $D_{\text{KS}} = \sup_x |F_1(x) - F_2(x)|$, compare CDFs directly, and information-theoretic divergences such as Kullback--Leibler (KL) divergence take a density-based approach.

% Paragraph reviewed and approved by Tak 2026-09-26 (Option C of the v6 paragraph review, queue #1): constant described as a per-unit clinical slope of the kind interaction analyses estimate (consistent with Section 2.4); (iii) aligned; "Crucially" removed.
% Paragraph revised 2026-10-11 (D-20261011-11, cold referee M1; Tak's standing instruction of 2026-10-11 23:49, individual veto open until submission): the KS bound's constant described, not dismissed
% Cut 2026-10-14 (D-20261014-1, needs Tak's veto): 'However, each of these approaches has limitations when applied to ... SMD captures only location differences ... KL divergence is asymmetric---comparing region A to region B yields a different value than comparing B to A---can diverge ... These gaps motivate ...(iii) offers such a bound. The Wasserstein-1 distance ..., developed in the following subsections, satisfy all three requirements.' -> same content, 'KL asymmetric' explanation clause and connecting words dropped, last two sentences joined with a semicolon
% Cut 2026-10-15 (D-20261015-1, needs Tak's veto): 'SMD captures only location and is blind to variance and shape.' (third statement of the same point), 'can diverge to infinity when empirical distributions do not fully overlap' -> 'infinite for non-overlapping empirical distributions', the separate sentence 'None of these measures offers a bound ... per-unit clinical slope' folded into requirement (iii) and 'the kind of quantity that treatment-by-covariate interaction analyses estimate' dropped (Sec. 2.4 says only that meta-regression on interaction estimates the average interaction slope); the KS total-variation constant sentence and requirements (i)-(iii) kept
Each of these approaches has limitations for comparing effect modifier distributions across regions. EDF-based statistics such as the KS statistic capture full distributional differences but lack a clinically interpretable scale, and the bound on treatment effect heterogeneity they admit is governed by the total variation of the CATE, a constant that is clinically stated only for a monotone CATE, for which $\mathrm{TV}(\tau) = |\tau(b) - \tau(a)|$ on a support $[a, b]$. KL divergence is asymmetric, infinite for non-overlapping empirical distributions, and requires density estimation that is impractical with small regional samples. These gaps motivate a measure that (i) captures distributional features beyond location, (ii) provides a clinically interpretable scale, and (iii) offers a bound on regional treatment effect differences whose constant is a per-unit clinical slope; the Wasserstein-1 distance and its clinical calibration, developed next, satisfy all three.

\subsection{Per-Effect-Modifier Wasserstein-1 Distance}\label{sec:nabcd_metric}

% Paragraph revised 2026-10-11 (D-20261011-11, cold referee M1; Tak's standing instruction of 2026-10-11 23:49, individual veto open until submission): adoption justified by attainment, not by the constant alone
% Paragraph revised 2026-10-12 (D-20261012-2, Louis M3; individual veto open until submission): attainment scoped to the Lipschitz class (the KS bound is attained within its own class)
% Paragraph revised 2026-10-12 (D-20261012-3, REF M1; individual veto open until submission): attainment no longer the reason for adoption (it is Kantorovich-Rubinstein duality restated); three reasons stated
% Cut 2026-10-15 (D-20261015-1, needs Tak's veto): '(sometimes called the Earth Mover's Distance)' and 'each highlighting a different facet of its interpretation' dropped (pure deletion of the two phrases); the three reasons for adopting W1 unchanged -> deleted (no replacement wording)
Among distributional distances that could capture the full distributional differences relevant to regional pooling, we adopt the Wasserstein-1 distance ($W_1$) for three reasons: it responds to differences beyond location; the heterogeneity bound it admits (Proposition~\ref{prop:heterogeneity} below) has a per-unit clinical slope $L_{\text{clinical}}$ as its constant; and that bound is valid for every CATE in the class, sharp when the slope is the only restriction placed on the CATE, and tighter, as a function of $W_1$ and the mean difference, when the CATE is also monotone; for any particular CATE it is therefore a conservative upper bound. The $W_1$ distance admits three equivalent representations on the real line.\cite{villani2009,vallender1974}

Let $F_1$ and $F_2$ denote the cumulative distribution functions (CDFs) of an effect modifier in two regions. Then
\begin{equation}
W_1(F_1, F_2) = \int_{-\infty}^{\infty} |F_1(x) - F_2(x)| \, dx = \int_0^1 |F_1^{-1}(u) - F_2^{-1}(u)| \, du = \inf_{\pi \in \Pi(F_1, F_2)} \int |x - y| \, d\pi(x, y),
\label{eq:wasserstein}
\end{equation}
% Cut 2026-10-15 (D-20261015-1, needs Tak's veto): 'The three forms are mathematically equivalent.' (stated in the paragraph before) dropped; the three form descriptions otherwise unchanged -> deleted (no replacement wording)
where $F_r^{-1}$ is the quantile function and $\Pi(F_1, F_2)$ is the set of probability couplings with marginals $F_1$ and $F_2$. The CDF form expresses $W_1$ as the area between the two CDFs (Figure~\ref{fig:nabcd_definition}); the quantile form as the integrated absolute difference of quantile functions; and the Kantorovich form as the minimum expected absolute displacement to transport one distribution onto the other.\cite{vallender1974} Unlike the standardized mean difference, $W_1$ responds to differences in variance, skewness, and other distributional features.\cite{panaretos2019}

\begin{figure}[t]
\centerline{\includegraphics[width=0.6\textwidth]{fig1_w1_definition.pdf}}% Cut 2026-10-15 (D-20261015-1, needs Tak's veto): figure width 0.8 -> 0.6 of the text width (layout only)
\caption{The Wasserstein-1 distance as the area between cumulative distribution functions, illustrated with a right-skewed Region~1 (gamma distribution) and a symmetric Region~2 (normal distribution). Panel~(A): probability densities. Panel~(B): cumulative distribution functions; the shaded area equals $W_1(F_1, F_2)$.\label{fig:nabcd_definition}}
\end{figure}

% Paragraph reviewed and approved by Tak 2026-09-26 (Option B of the v6 paragraph review, queue #2): units sentence corrected (the CDF-form integrand is dimensionless; W1 carries units in all three forms); transport-distance wording from the ERRFIX kept.
% Cut 2026-10-14 (D-20261014-1, needs Tak's veto): 'A property of W1 central to our framework is that it carries the units of the underlying effect modifier. In each of the three equivalent forms ..., W1 carries the units of X: ... Consequently, W1 is expressed on the original scale ... This dimensional alignment supports ... we therefore restrict attention in this paper to ...' -> opening and 'Consequently' sentences merged into one, same units statement, example and scope restriction
% Cut 2026-10-15 (D-20261015-1, needs Tak's veto): 'In each of the three equivalent forms: the CDF form integrates a dimensionless difference of probabilities over x, the quantile form differences of quantiles, and the coupling form distances |x-y|' -> 'in each of the three equivalent forms'; 'so W_1(F_1, F_2) is expressed on the original scale of the effect modifier, which supports direct clinical interpretation' -> 'so it is expressed on the original scale and supports direct clinical interpretation'; 'The same logic precludes a meaningful direct comparison' -> 'are not directly comparable'; 'Section 5 elaborates on this scope choice' dropped; 5 mmHg example unchanged
$W_1$ carries the units of the underlying effect modifier in each of the three equivalent forms of equation~(\ref{eq:wasserstein}), so it is expressed on the original scale and supports direct clinical interpretation: a regional $W_1 = 5$~mmHg in systolic blood pressure means that optimally transporting one distribution onto the other moves probability mass an average distance of 5~mmHg. Values of $W_1$ for effect modifiers measured on different scales are not directly comparable, so we restrict attention to per-effect-modifier assessment.

% Cut 2026-10-15 (D-20261015-1, needs Tak's veto): 'satisfying non-negativity, symmetry, identity of indiscernibles, and the triangle inequality' (list of metric axioms) -> dropped (pure deletion; the non-negativity / identity proof of Appendix A.1 is removed with it); the affine-equivariance property 'W_1(F_{aX+b}, F_{aY+b}) = |a| W_1(F_X, F_Y)' dropped (not used later); 'Beyond these basic properties' -> 'also'
$W_1$ is a proper metric on the space of probability distributions with finite first moment.\cite{villani2009} It also admits a Kantorovich--Rubinstein dual representation\cite{villani2009}
\begin{equation}
W_1(F_1, F_2) = \sup_{\|f\|_{\text{Lip}} \leq 1} \left|\int f \, dF_1 - \int f \, dF_2\right|,
\label{eq:dual}
\end{equation}
% Cut 2026-10-15 (D-20261015-1, needs Tak's veto): 'for all x, y' after the Lipschitz inequality of tau dropped; the notation parenthetical kept -> deleted (no replacement wording)
which expresses $W_1$ as the largest possible difference in expected values of a 1-Lipschitz function $f$ between the two distributions ($\|f\|_{\text{Lip}} \leq 1$ means $|f(x) - f(y)| \leq |x - y|$ for all $x, y$). If the conditional average treatment effect function $\tau(x)$ is Lipschitz continuous with constant $L_{\text{clinical}}$, so that $|\tau(x) - \tau(y)| \leq L_{\text{clinical}} |x - y|$, then $\tau / L_{\text{clinical}}$ is 1-Lipschitz and applying equation~(\ref{eq:dual}) yields the central result of this section:

% Proposition reviewed and approved by Tak 2026-09-26 (Option B of the v6 paragraph review, queue #3): assumption that the CATE is the same in every region and depends on X alone.
\begin{proposition}\label{prop:heterogeneity}
\textup{(Connection to heterogeneity).} Suppose the CATE function $\tau$ is the same in every region and depends on the effect modifier $X$ alone. If $\tau$ is Lipschitz continuous with constant $L_{\text{clinical}}$, then for any two regions with effect modifier distributions $F_1$ and $F_2$,
\begin{equation}
|\bar{\tau}_1 - \bar{\tau}_2| \leq L_{\text{clinical}} \cdot W_1(F_1, F_2).
\label{eq:heterogeneity_bound}
\end{equation}
\end{proposition}

\begin{proof}
Apply equation~(\ref{eq:dual}) with $f = \tau / L_{\text{clinical}}$, which is 1-Lipschitz by assumption. This yields $\left|\int (\tau / L_{\text{clinical}}) \, d(F_1 - F_2)\right| \leq W_1(F_1, F_2)$. Combining with equation~(\ref{eq:regional_ate}) for $\bar{\tau}_r$ and rearranging yields equation~(\ref{eq:heterogeneity_bound}).
\end{proof}

% Paragraph reviewed and approved by Tak 2026-09-21 (Option B of the paragraph
% review; Claude Science referee item M1): the claim that the bound is specific
% to W_1 was false (integration by parts gives a KS bound with constant TV(tau);
% W_1 <= W_2 gives a looser W_2 bound). Replaced by the constant / tightness
% argument with the monotone-CATE case stated honestly.
% Paragraphs reviewed and approved by Tak 2026-09-26 (Option C of the v6 paragraph review, queue #4): KS/W1 ordering (neither dominates), "tightness" removed, constant described as the kind of quantity interaction analyses estimate, normalization no longer attributed to prior literature; split into two paragraphs.
% Paragraph revised 2026-10-11 (D-20261011-11, cold referee M1; Tak's standing instruction of 2026-10-11 23:49, individual veto open until submission): attainment stated (tau* construction, |tau| <= 1 admissibility), constant second
% Paragraph revised 2026-10-12 (D-20261012-2, Louis M3/m3; individual veto open until submission): attainment scoped to the class; the KS bound's own attainment stated; x_0 fixed at the midpoint so that |tau*| <= 1 follows from L(b-a) <= 2
% Cut 2026-10-14 (D-20261014-1, needs Tak's veto): the Kolmogorov--Smirnov bound is likewise sharp within its own class, by a step CATE of height TV(tau) at the point where |F1 - F2| is largest -> ... within its own class (Appendix A.3) [the step-CATE sentence is moved to A.3 verbatim]
% Cut 2026-10-15 (D-20261015-1, needs Tak's veto): the two paragraphs on exclusivity and sharpness -> one block. Deleted: 'Sharpness, which the Kolmogorov--Smirnov bound also has within its own class, does not distinguish the two bounds; the constant does ...', the KS-sharpness clause, 'Sharpness over a class is a property of the class, and shape information available to the sponsor tightens the bound' (joined to the monotone sentence), the L_UB attainment sentence ('A bound on the magnitude of tau itself tightens it further ...'; it is in Appendix A.2). Kept: tau* sharpness with Kantorovich--Rubinstein, eq:mono_bound, 'conservative by construction', 'neither dominates (Appendix A.2)'. New sentence J (insurance, Louis): the bound is not specific to W1 (TV.D_KS, L.W2), each sharp in its own class, the constants differ
\noindent % Cut 2026-10-15 (D-20261015-1, Louis review-fix, needs Tak's veto): sentence J: 'each is sharp within its own class, and the bounds differ in their constant' (false for L*W2, same constant) -> bounded-variation condition, KS bound sharp in its class, L*W2 valid but at least L*W1
The bound is not specific to $W_1$: for a CATE of bounded variation on a bounded support, integration by parts gives a bound $\mathrm{TV}(\tau) \cdot D_{\text{KS}}$, sharp within its own class, and $W_1 \leq W_2$ gives a valid but looser $L_{\text{clinical}} \cdot W_2$ with the same constant, which, like every valid bound with that constant, is at least $L_{\text{clinical}} W_1$ (Appendix~\ref{app:ks_bound}).\cite{villani2009} Equation~(\ref{eq:heterogeneity_bound}) is sharp when the Lipschitz constant is the only restriction placed on the CATE: by the Kantorovich--Rubinstein duality of equation~(\ref{eq:dual}), $L_{\text{clinical}} W_1$ is the supremum of $|\bar{\tau}_1 - \bar{\tau}_2|$ over all common $L_{\text{clinical}}$-Lipschitz CATEs, attained by $\tau^*(x) = L_{\text{clinical}} \int_{x_0}^{x} \operatorname{sgn}(F_2(t) - F_1(t))\,dt$, whose slope changes sign at every crossing of the two CDFs. If $\tau$ is also monotone, with $0 \leq \tau' \leq L_{\text{clinical}}$, then $\bar{\tau}_1 - \bar{\tau}_2 = \int (F_2 - F_1)\,\tau'\,dx$ lies between $-L_{\text{clinical}} A_-$ and $L_{\text{clinical}} A_+$, where $A_\pm = \int (F_2 - F_1)_\pm\,dx$ satisfy $A_+ + A_- = W_1$ and $A_+ - A_- = \mu_1 - \mu_2$, so that
\begin{equation}
\sup_{\tau\ \text{monotone},\ L_{\text{clinical}}\text{-Lipschitz}} |\bar{\tau}_1 - \bar{\tau}_2| = L_{\text{clinical}} \max(A_+, A_-) = \tfrac{1}{2} L_{\text{clinical}} \left( W_1 + |\mu_1 - \mu_2| \right),
\label{eq:mono_bound}
\end{equation}
which equals $L_{\text{clinical}} W_1$ when the CDFs do not cross and $L_{\text{clinical}} W_1 / 2$ when the means coincide. $\Delta_{\max} = L_{\text{clinical}} W_1$ is therefore a conservative upper bound by construction; Section~\ref{sec:app_clinical} reports the factor by which equation~(\ref{eq:mono_bound}) would shrink it in the application. Neither the $W_1$ bound nor the Kolmogorov--Smirnov bound dominates the other in general (Appendix~\ref{app:ks_bound}).
% Restored 2026-10-14 01:30 (Louis review-fix of the cut, K1/K2): the M12 cut removed a 2026-10-12 correction sentence (rule 2); one clause restored.
% Paragraph revised 2026-10-12 (D-20261012-3, REF M1 / EXT Minor 2 / EXT Minor 5; individual veto open until submission): sharpness stated as a property of the class (KR duality), the monotone sharp bound L(W1+|dmu|)/2 derived, the bounded-effect and unbounded-support caveats added, the KS-bound comparison moved to Appendix A.3

% Cut 2026-10-14 (D-20261014-1 M06, R2c): CLAIMS C-04
% Cut 2026-10-15 (D-20261015-1, needs Tak's veto): 'so no further normalization is required for a clinically interpretable bound' and 'normalizing by a within-effect-modifier scale measure such as the pooled interquartile range is redundant once the bound is used for calibration (Section 2.4)' -> one clause; the Lipschitz-appropriateness sentence unchanged
Because the bound involves $W_1$ in its original measurement units multiplied by a clinical slope, no further normalization is required and we adopt $W_1$ in its original units throughout. The Lipschitz assumption of Proposition~\ref{prop:heterogeneity} is appropriate when the CATE varies smoothly with the effect modifier; for threshold or step-function CATEs, $L_{\text{clinical}}$ can be arbitrarily large and the bound becomes conservative (Section~\ref{sec:discussion}).

\subsection{Estimation}\label{sec:estimation}

% Citation-only correction 2026-09-27 (Louis, v6 review of queue #48-49): the real-line bootstrap statements are sourced to Fang and Santos 2019 and Dumbgen 1993; Sommerfeld and Munk 2018 covers finite supports only and is kept for that case.
% Paragraph reviewed and approved by Tak 2026-09-26 (Option C of the v6 paragraph review, queue #5): tie example corrected to SBP in whole mmHg (GUSTO age is recorded to three decimals); bootstrap-consistency condition split into two sentences.
Given samples $\{X_{1,i}\}_{i=1}^{n_1}$ from region 1 and $\{X_{2,j}\}_{j=1}^{n_2}$ from region 2, the Wasserstein-1 distance is estimated using empirical distribution functions. Since $\hat{F}_1$ and $\hat{F}_2$ are piecewise constant (right-continuous step functions), the integral in equation~(\ref{eq:wasserstein}) reduces to a finite sum over the intervals between consecutive distinct values of the pooled sample:
\begin{equation}
\widehat{W}_1 = \sum_{k=1}^{m-1} |\hat{F}_1(x_{(k)}) - \hat{F}_2(x_{(k)})| \cdot (x_{(k+1)} - x_{(k)}),
\label{eq:estimator}
\end{equation}
% Paragraph revised 2026-10-11 (D-20261011-11, cold referee Minor 2; Tak's standing instruction of 2026-10-11 23:49, individual veto open until submission): why the bootstrap rather than the available Wald interval
% Cut 2026-10-15 (D-20261015-1, needs Tak's veto): Brownian-bridge limit statement and its integrability condition, the Hadamard-derivative sentence, 'by the functional delta method', the convergence-rate clause, 'ties are pervasive for coarsely recorded modifiers', 'The limit distribution of W1hat depends on the unknown F, so no universal critical values are available' (moved to Appendix A.1 in one sentence) and the O(n log n) sentence -> dropped here (Appendix A.1 states the limit and the conditions; its grid paragraph covers ties); the convergence-rate clause is dropped everywhere; 'Each bootstrap replicate draws n_1 observations ... from region 1 and, independently, n_2 from region 2, and recomputes W1hat by equation (estimator)' -> 'Each replicate resamples each region independently, with replacement and at its own size'; 'the same condition governs the consistency of the bootstrap, so the choice between them rests on finite-sample behaviour' -> 'A Wald interval is governed by the same condition and requires an estimate of the limit's variance ...'; the measure-zero condition for the Gaussian limit / bootstrap consistency is stated before the Wald-versus-bootstrap choice, which keeps its one-sentence reason; B = 2,000, independent per-region resampling and type-7 quantiles kept
where $\hat{F}_r$ denotes the empirical CDF of region $r$ and $x_{(1)} < \cdots < x_{(m)}$ are the $m$ distinct values of the pooled sample. When $\{x : F_1(x) = F_2(x) \in (0,1)\}$ has Lebesgue measure zero, as when the two CDFs cross only at isolated points (Appendix~\ref{app:asymptotic}, with its moment and sample-ratio conditions), the limit is Gaussian and the ordinary nonparametric bootstrap is consistent;\cite{fang2019} where the CDFs coincide on an interval, including $F_1 = F_2$, the naive bootstrap is not consistent.\cite{dumbgen1993,fang2019,sommerfeld2018} A Wald interval is governed by the same condition and requires an estimate of the limit's variance, which depends on $\operatorname{sgn}(F_1 - F_2)$ and hence on where the CDFs cross, whereas the percentile bootstrap requires none; we therefore employ the nonparametric percentile bootstrap with $B = 2{,}000$ replicates.\cite{efron1993} Each replicate resamples each region independently, with replacement and at its own size, and the 95\% interval is formed by the 2.5th and 97.5th sample percentiles of the $B$ replicates (R's default quantile definition, type 7).

% Paragraph reviewed and approved by Tak 2026-09-26 (Option B of the v6 paragraph review, queue #6): null floor as an estimator property; source distribution stated (anchor in Section 4.2, combined pair in Supplement C); 95th percentile as reading reference.
% Paragraph revised 2026-10-11 (D-20261011-11, cold referee M4; Tak's standing instruction of 2026-10-11 23:49, individual veto open until submission): the null floor named as a one-sided level-0.05 test that governs display
% Cut 2026-10-14 (D-20261014-1, needs Tak's veto): 'it is biased upward at every sample size: since E F_r(x) = F_r(x), Jensen's inequality gives E|F1hat - F2hat| >= |F1 - F2| at every x, and integrating (the integrand is non-negative, so no further condition is needed) E W1hat >= W1 for all n1, n2; the percentile interval does not correct this bias, ...' -> 'biased upward at every sample size, E W1hat >= W1 (Appendix A.1), and the percentile interval does not correct this bias, ...' [the pointwise-Jensen / non-negative-integrand sentence is moved verbatim to Appendix A.1; the enrolled-sample sentence is unchanged] (Harvey menu 6)
% Cut 2026-10-15 (D-20261015-1, needs Tak's veto): pointer 'Appendix A.1' (Jensen derivation) -> app:asymptotic (the Jensen paragraph moved to the first appendix subsection); 'More precisely' sentence joined to the first by a colon; 'so near the boundary the whole interval is displaced upward (Section 3.2.2)', 'whatever the sampling error of W1hat as an estimate of W1', 'the reference point used in Section 4.2', 'because it depends on both sample sizes and on the source distribution', 'equivalently of Delta_max' and 'used only to mark which estimates are read at their upper limit' compressed or dropped; enrolled-sample sentence, multiplicity sentence and the one-sided level-0.05 test kept
Because $W_1$ is a distance, $\widehat{W}_1$ is non-negative by construction and typically positive when the two regions are draws from the same distribution: it is biased upward at every sample size, $E\,\widehat{W}_1 \geq W_1$ (Appendix~\ref{app:asymptotic}), and the percentile interval does not correct this bias. For the patients actually enrolled, the bound holds exactly with the empirical distributions: the difference between the two enrolled samples' average treatment effects is at most $L_{\text{clinical}}\,W_1(\hat{F}_1, \hat{F}_2)$. We refer to the sampling distribution of $\widehat{W}_1$ under $F_1 = F_2$ as the null floor; it is obtained, per region pair, by drawing two independent resamples, of the two regions' actual sizes, from one region's empirical distribution---the anchor's in Section~\ref{sec:app_operability}---and recomputing $\widehat{W}_1$. An observed $\widehat{W}_1$ at or below its 95th percentile is one the estimator would produce with appreciable probability from identical distributions, so it is not evidence of a difference; how large the difference could be is read from the upper confidence limit of $\widehat{W}_1$. Declaring a distance unresolved in this way is a one-sided test of $F_1 = F_2$ at level 0.05; the comparisons of Section~\ref{sec:application} are not adjusted for multiplicity, since no inference about pooling is drawn from the number of unresolved distances.
% Paragraph revised 2026-10-12 (D-20261012-3, EXT M4(a) / EXT Minor 1 / REF Minor 4): Jensen bias stated as a result; finite-sample bound for the enrolled samples; multiplicity stated.

\subsection{Interpretation and Clinical Calibration}\label{sec:inference}

% Cut 2026-10-14 (D-20261014-1 M09, F): CLAIMS C-13 (no cutoffs is stated in the next paragraph and in Sec. 2.4 last paragraph)

% Paragraph reviewed and approved by Tak 2026-09-06 (Option B of the paragraph
% review; simulated-review point R2.1(b)): "provides a per-unit slope directly"
% replaced -- an average interaction slope is a lower bound on the Lipschitz
% constant, not the constant; safety-factor and premise sentences added;
% "significance" removed (estimation-centered wording).
% Paragraph reviewed and approved by Tak 2026-09-26 (Option B of the v6 paragraph review, queue #7): unverified VanderWeele clause removed; average slope is a lower bound in absolute value.
% Cut 2026-10-15 (D-20261015-1, needs Tak's veto): 'By definition, L_clinical quantifies ...' and the two sentences on the average slope (lower bound; safety factor 'so that the premise of Proposition 1 remains defensible') merged into one ('rather than the constant itself' dropped); the sentence 'More commonly at the MRCT planning stage, published subgroup reports document the direction of interaction rather than a per-unit slope, and a precise numerical L_clinical is unavailable. We address both situations below.' dropped (the next paragraph treats the two cases: L_clinical known, L* when unknown); the units sentence shortened -> deleted (no replacement wording)
We refer to the Lipschitz constant of the CATE function $\tau(x)$ as the clinical slope $L_{\text{clinical}}$: the maximum change in the treatment effect per unit change in the effect modifier, expressed on the outcome scale of the trial---for example, absolute 30-day mortality risk per year of age, or per mmHg of systolic blood pressure. It is quoted on the proportion scale of the outcome, so that $L_{\text{clinical}} = 1 \times 10^{-2}$ per year denotes one percentage point of absolute 30-day mortality per year of age, and clinical margins $\Delta_{\text{clin}}$ stated in percentage points are divided by 100 before entering any calculation. Meta-regression on treatment-by-covariate interaction in prior individual patient data (IPD) estimates the average interaction slope,\cite{riley2010,fisher2017} which is, in absolute value, a lower bound on the Lipschitz constant, so an $L_{\text{clinical}}$ specified from it should carry a sponsor-specified multiplicative safety factor, in the spirit of margin discounting in non-inferiority design.

From Proposition~\ref{prop:heterogeneity}, the maximum potential difference in regional average treatment effects attributable to effect modifier distributional differences is
\begin{equation}
\Delta_{\max} = L_{\text{clinical}} \cdot W_1(F_1, F_2),
\label{eq:delta_max}
\end{equation}
% Cut 2026-10-15 (D-20261015-1, needs Tak's veto): 'Any normalized form W_1 / X multiplied by the same scale X cancels in equation (delta_max), so W_1 in its original units yields the same clinical-scale calibration as any rescaled version.' (repeats the paragraph after Proposition 1) -> dropped; the last sentence joined to the equation sentence
expressed in the units of the trial outcome scale; when $L_{\text{clinical}}$ is independently known, equation~(\ref{eq:delta_max}) provides a direct clinical-scale translation of the observed distributional difference.

% Paragraph reviewed 2026-09-27 (Option C, Harvey's recommendation under Tak's standing instruction; v6 paragraph review, queue #8): L* interval defined here once; "full range of evidence states" narrowed to the two evidence states.
The estimator $\widehat{W}_1$ has a percentile bootstrap CI (Section~\ref{sec:estimation}), and inferential uncertainty propagates linearly to $\Delta_{\max}$ when $L_{\text{clinical}}$ is treated as fixed: $[\Delta_{\max,L}, \Delta_{\max,U}] = L_{\text{clinical}} \cdot [W_{1,L}, W_{1,U}]$. When $L_{\text{clinical}}$ itself is uncertain or unavailable, we solve equation~(\ref{eq:delta_max}) for $L_{\text{clinical}}$ at a pre-specified clinical margin $\Delta_{\text{clin}}$ (the minimal clinically important treatment effect difference, the non-inferiority margin, or the targeted overall effect), yielding the required clinical slope
\begin{equation}
L^* = \frac{\Delta_{\text{clin}}}{W_1(F_1, F_2)}.
\label{eq:lstar}
\end{equation}
% Cut 2026-10-14 (D-20261014-1 M10, F): CLAIMS C-06
% Cut 2026-10-15 (D-20261015-1, needs Tak's veto): 'reading L_UB against the lower limit Delta_clin/W_{1,U} rather than the point value carries the estimation uncertainty of W1hat into this comparison' dropped (pure deletion of the clause) -> deleted (no replacement wording)
$L^*$ is the smallest clinical slope under which the observed distributional difference could produce a clinically meaningful regional treatment effect difference, within the class of CATEs restricted only by their slope. Because $L^*$ is a decreasing function of $W_1$, the bootstrap CI of $\widehat{W}_1$ gives the percentile interval $[\Delta_{\text{clin}}/W_{1,U},\, \Delta_{\text{clin}}/W_{1,L}]$ for $L^*$. If a plausible upper bound on $L_{\text{clinical}}$ from prior evidence (denoted $L_{\text{UB}}$) lies below $L^*$, the distributional difference cannot produce such a difference, provided $L_{\text{clinical}} \leq L_{\text{UB}}$.

% Paragraph reviewed and approved by Tak 2026-09-06 (Option C of the paragraph
% review; simulated-review point R1.2(a)): the "so that any plausible CATE
% sensitivity falls below them" deduction retracted; L_UB repositioned as a
% sponsor-specified clinical judgment with the arm-decomposition information
% relation, the Section 4.4 sensitivity analysis promoted to the primary
% robustness statement, and the prognostic-vs-predictive taxonomy cited.
% Paragraph reviewed 2026-09-27 (Option C, Harvey's recommendation under Tak's standing instruction; v6 paragraph review, queue #9): prognostic-gradient reference point kept but qualified against the arm-sum inequality; closing "accommodates the typical uncertainty" sentence removed (duplicated the preceding sentence).
% Cut 2026-10-14 (D-20261014-1, needs Tak's veto): '...which is one reason the bounds are illustrative and L* is reported for every pair. Conclusions drawn from the bounds are conditional on them; because L* is reported for every region pair, any alternative bound can be compared against it directly. Alternative specification strategies include the deft estimator of Fisher et al.' -> the repeated 'L* is reported for every pair' clause and the last sentence dropped (fisher2017 stays cited in the paragraph on the average slope)
% Cut 2026-10-14 (D-20261014-1, needs Tak's veto): 'these exceed the prognostic age gradient implied by the 30-day mortality that Lee et al. report for GUSTO-I, 1.1% below 45 years and 20.5% above 75 years, a rise of 19.4 percentage points over at least three decades and hence an average gradient of at most about 6.5e-3 per year, though not the sum of two arm-level gradients ...' -> 'these exceed the average prognostic age gradient implied by the 30-day mortality that Lee et al. report for GUSTO-I (at most about 6.5e-3 per year; Appendix A.3), though not the sum of two arm-level gradients ...' [the 1.1%/20.5%/19.4-point arithmetic moves to a sentence in Appendix A.3; the 'reference point rather than a guarantee' qualifier stays] (Harvey menu 7)
% Cut 2026-10-15 (D-20261015-1, needs Tak's veto): arm-level decomposition ('Writing R_a(x) for the risk under arm a, tau = R_1 - R_0, so |tau'| <= |R_1'| + |R_0'|: arm-level risk gradients constrain the modification slope only jointly') and 'not the sum of two arm-level gradients of that size that the inequality above permits' dropped with it; the GUSTO-I overall treatment effect sentence (6.3% under accelerated t-PA against 7.2% and 7.4%, 'which is one reason the bounds are illustrative'), 'A Delta_max larger than the overall treatment effect does not exclude a regional difference of that size' and 'Because L* is reported for every region pair, any alternative bound can be compared against it directly' dropped, no remaining carrier in the paper (its Discussion carrier was removed on 2026-10-14); prognostic/predictive distinction (Ballman, Kent), the L_UB reference points (FTT, Lee et al.) and the age-gradient comparison kept -> deleted (no replacement wording)
Specifying $L_{\text{clinical}}$ (or its upper bound $L_{\text{UB}}$) is a clinical judgment, informed but not determined by prior evidence, for which the sponsor carries responsibility---the same footing as a non-inferiority margin. A prognostic gradient alone guarantees no bound on the modification slope, prognostic and treatment-effect gradients being distinct quantities.\cite{ballman2015,kent2020} For the GUSTO-I case study presented in Section~\ref{sec:application}, the illustrative bounds are $L_{\text{UB,age}} = 1 \times 10^{-2}$ per year and $L_{\text{UB,SBP}} = 2 \times 10^{-3}$ per mmHg on the 30-day mortality scale, set so that one year of age or 5~mmHg of systolic blood pressure would reach a margin of 1 percentage point; the Fibrinolytic Therapy Trialists' meta-analysis\cite{ftt1994} and the GUSTO-I prognostic model\cite{lee1995} are reference points, not sources of the bounds. The age bound exceeds the average prognostic age gradient implied by the 30-day mortality that Lee et al.\ report for GUSTO-I (at most about $6.5 \times 10^{-3}$ per year; Appendix~\ref{app:ks_bound}),\cite{lee1995} a reference point rather than a guarantee.

% Reopened 2026-09-27 (Tak: select pooling regions by ranking; simulation identifies the method, GUSTO-I applies it): ordering of candidates by W1 restored.
% Paragraph reviewed 2026-09-27 (Option B, Harvey's recommendation under Tak's standing instruction; v6 paragraph review, queue #10): "ranking partners" replaced by reading each distance against its null floor and upper confidence limit; "appropriate threshold" (threshold residue) removed.
% Paragraph revised 2026-10-11 (D-20261011-11, cold referee M4; Tak's standing instruction of 2026-10-11 23:49, individual veto open until submission): estimation-centered statement reconciled with the null-floor test
% Cut 2026-10-14 (D-20261014-1 M11, F): CLAIMS C-13, C-42
% Cut 2026-10-15 (D-20261015-1, needs Tak's veto): the first sentence ('The pooling assessment proceeds by computing per-effect-modifier W1 with bootstrap CIs across candidate partner regions, ordering the candidates by the point estimate within each effect modifier, and reading each distance against its null floor and its upper confidence limit (distances unresolved against the null floor are not ordered among themselves)', a restatement of Secs. 2.3 and 4.2), '(the overall treatment effect size, the non-inferiority margin, the minimal clinically important difference)' and 'presented alongside the clinical context' dropped; 'estimation-centered rather than testing-based', 'no universal cutoff' and 'the one test, the null floor ...' unchanged -> deleted (no replacement wording)
The framework is deliberately estimation-centered rather than testing-based: $W_1$ has no universal cutoff that separates ``similar'' from ``dissimilar,'' and what a given $W_1$ means clinically depends on the clinical slope $L_{\text{clinical}}$, which is external to the distributional comparison. Per-pair $\Delta_{\max}$ (or $L^*$) with bootstrap CIs are inputs for sponsor and regulatory judgment, which is not driven by a significance test;\cite{wasserstein2016} the one test, the null floor of Section~\ref{sec:estimation}, governs how an estimate is displayed and read, not whether regions are pooled.

%==============================================================================
\section{Simulation Studies}\label{sec:simulation}
% Labels follow the text (sec:study1 = identification, sec:study2 = estimation; fixed 2026-10-11). tab:study2_*, fig:study2_auc and R/verify_study2_figures.R keep the pre-D-20260927-3 numbering.
%==============================================================================

% Reopened 2026-09-27 (Tak: the simulation identifies which method correctly identifies pooling regions by ranking for an anchor, and GUSTO-I applies it; clustering removed the same day, PROTOCOL_silhouette_extension.md criterion not met).
% Paragraph reviewed 2026-09-27 (Option B, Harvey's recommendation under Tak's standing instruction; v6 paragraph review, queue #11): identification stated without an anchor (the clustering task has none); required-n statements put on one measure-specific convention; decision wording removed.
% Paragraph revised 2026-10-12 (D-20261012-3, REF M4 / EXT M3; individual veto open until submission): Study 1 is identification on a selection task with the number of matches supplied; it does not evaluate the ordering of partners that all differ, nor ranking by effect difference, which depends on the CATE (linear vs V-shaped example)
% Cut 2026-10-15 (D-20261015-1, needs Tak's veto): 'two complementary' dropped; the scale-only example inside S5 (L sqrt(2/pi) . 5), 'a region differing from the anchor in scale alone has an effect difference of zero although W_1 > 0' under the linear CATE, 'The Lipschitz class contains both CATEs, Proposition 1 covers both, and the SMD covers only the first', 'Study 1 carries the comparison between measures, because a distance that cannot distinguish differing distributions cannot usefully be calibrated' and 'which underwrites the conversion of an estimated distance into a Delta_max statement via equation (delta_max)' -> shortened or dropped; reworded: 'a difference in scale alone produces the full L W_1' -> 'a scale difference about the centre 50 produces the full L W_1' (the V-shaped CATE is centred at 50); the self-limitation (ordering of regions that all differ from the anchor is not evaluated) and the linear versus V-shaped CATE contrast kept
We conducted two simulation studies. Study 1 evaluates identification: whether each candidate distance measure can detect, at some sample size, that a candidate region's effect modifier distribution differs from the anchor's, scored on a selection task in which candidates built from the anchor's population are ranked against candidates built from a different population and the number of true matches is supplied. It does not evaluate the ordering of regions that all differ from the anchor, which is what the application reports, and it does not rank regions by their effect difference, which depends on the unknown CATE: under a linear CATE only the mean difference matters, whereas under the V-shaped CATE $\tau(x) = L_{\text{clinical}} |x - 50|$ a scale difference about the centre 50 produces the full $L_{\text{clinical}} W_1$ while the SMD is zero; Study 1 asks which measures are blind to differences of the second kind. Study 2 evaluates estimation performance: bias, variability, and coverage probability of bootstrap confidence intervals for $\widehat{W}_1$.

% Cut 2026-10-15 (D-20261015-1, needs Tak's veto): paragraph 'Study 1's required sample sizes are discrimination statements, the per-region sample size at which a measure's AUC reaches at least 0.9 ... Study 2's estimation statements are not interchangeable with these' deleted (pure deletion; the definition of the required n is in Sec. 3.1.3 and in the caption of its table; the AUC abbreviation is now defined where Study 1's performance measure is introduced) -> deleted (no replacement wording)

\subsection{Study 1: Identification}\label{sec:study1}

% Figures: R/selection_simulation.R + R/clustering_simulation.R
% (results/selection_sim_summary{,_ext}.csv, clustering_sim_summary{,_ext}.csv,
% required_n_{selection,clustering}{,_floor}.csv, clustering_nok_silhouette.csv).
% Every number quoted in this subsection is checked by R/verify_study2_figures.R
% against those CSVs and the generators (results/verify_study2.log, ALL PASS
% 2026-08-30; gate extended and re-run ALL PASS 2026-09-27 after the v6 paragraph review). Two recorded-note errors were corrected by that gate: the blind
% plateau values are max_auc (0.50-0.51), n=2000 values are 0.49-0.50; and the
% design constants are 10,000/3,000 reps (selection), 5,000/2,000 (clustering),
% 9 candidates in every set (not "5,000-10,000" / "9-11").

\subsubsection{Design}\label{sec:study1_design}

% Paragraph reviewed 2026-09-27 (Option A, Harvey's recommendation under Tak's standing instruction; v6 paragraph review, queue #12): roster corrected to two bulk-shift control regions; "target of the decision" reworded.
% Paragraph revised 2026-10-11 (D-20261011-11, cold referee M3; Tak's standing instruction of 2026-10-11 23:49, individual veto open until submission): the KS-saturation mechanism of Set 4 stated
% Paragraph revised 2026-10-12 (D-20261012-2, Louis C4; individual veto open until submission): "deliberately favourable to CDF-supremum statistics" removed — Sets 1 and 2 show a uniform shift does not favour the KS statistic in general
% Cut 2026-10-12 (D-20261012-3 length plan S02): gloss on the Set 4 geometry sentence (pattern of a small severe subpopulation; control wording shortened)
% Cut 2026-10-15 (D-20261015-1, needs Tak's veto): 'Set 1 (Gaussian) places discordance in location and scale. Set 2 ... Set 3 ... every region, matched or discordant, is constructed with mean 50 and SD 10, so all discordance lies in distributional shape' -> three sentences joined; 'whatever the sample size', 'each built from a single distributional family' and 'mass added 40 units from the centre moves W1 by that distance per unit of added mass' dropped; the Set 4 mechanism condensed to one sentence ('a subgroup of prevalence eps displaced by d changes KS by at most about eps but W1 by eps d, so W1's advantage follows from the construction; the bulk-shift control removes it'); the construction-label / not-circular sentence kept
Study 1 follows the ADEMP structure (aims, data-generating mechanisms, estimands, methods, performance measures) for simulation studies.\cite{morris2019} Four worlds are simulated, each with regions constructed either from the anchor's own population (true matches) or from a discordant one. Set 1 (Gaussian) places discordance in location and scale; Set 2 (log-normal) in location and in dispersion combined with shape; Set 3 (mixture) in distributional shape alone, every region, matched or discordant, having mean 50 and SD 10. Set 4 (extremes) places a 5\% subgroup in each tail, 40 units from the centre in the anchor; discordant regions displace both subgroups (to 70 or 100 units), raise their prevalence to 12\% (mean fixed), or displace only the upper subgroup to 100 units, and two control regions instead shift their whole population by 2 and 3 units. In the three Set 4 cells in which a subgroup is displaced (symmetric and asymmetric severity) or its prevalence raised, a subgroup of prevalence $\varepsilon$ displaced by $d$ changes the KS statistic by at most about $\varepsilon$ whatever $d$ is, but $W_1$ by $\varepsilon d$, so $W_1$'s sample-size advantage there follows from the construction; the bulk-shift control, a location shift of the whole population, removes it. Ground truth is the construction label, so the target is defined without reference to any of the compared measures and the comparison is not circular.

% Paragraph reviewed 2026-09-27 (Option C, Harvey's recommendation under Tak's standing instruction; v6 paragraph review, queue #13): RV1 described as our application of the cited distance to one continuous modifier (Lasso weight a common scale factor); RV1/SMD agreement scoped to equal regional SDs; the log-scale SMD result reported.
% Rewritten 2026-09-27 (Tak): RV1-RV3 removed (Komiyama et al. introduce one example approach, not an evaluated method); three measures remain.
% Paragraph revised 2026-10-11 (D-20261011-11, cold referee M5; Tak's standing instruction of 2026-10-11 23:49, individual veto open until submission): methods and tasks run under the protocol but not reported are declared (ADEMP)
% Cut 2026-10-14 (D-20261014-1 S17, B): CLAIMS C-28 (only carrier removed: Harvey's claim decision D-20261014-1 (b), B-tier; C-40 wording already lists the log-scale SMD as not reported)
% Cut 2026-10-15 (D-20261015-1, needs Tak's veto): the list of analyses run under the registered protocol (three representative-value distances built from one to three regional moments, SMD on the log scale with its two AUC values, a harm measure, a silhouette clustering task, a threshold-rule study, a sensitivity analysis of an eligibility rule) -> one sentence (insurance B, Louis); the log-scale SMD AUC values (aucSMDlogShapeN2000, aucSMDlogCombinedN2000) are no longer reported
Three measures enter the comparison: $W_1$; the KS statistic, the most common comparison of full distributions; and SMD, the standard comparison of location. % Cut 2026-10-15 (D-20261015-1, Louis review-fix, needs Tak's veto): disclosure completed (false-pooling and precision measures, harm measure definition, threshold-rule study, eligibility-rule sensitivity analysis, repository pointer)
Other distances, tasks and performance measures run under the protocol (representative-value distances, a log-scale SMD, a harm measure based on the largest $W_1$ among the retained regions, the probability that the retained set includes a discordant region, a threshold-rule study, an eligibility-rule sensitivity analysis, and a clustering task whose registered criterion for estimating the number of groups was not met) are not reported; this paper compares distribution-level measures with the single-summary SMD and prescribes no grouping procedure; the outputs are in the repository.

% Paragraph reviewed 2026-09-27 (Option B, Harvey's recommendation under Tak's standing instruction; v6 paragraph review, queue #14): false-pooling definition corrected to the quantity the code computes (probability that the retained set includes a discordant region); "admits" replaced; oracle k acknowledged.
% Reopened 2026-09-27 (Tak): clustering task removed from the paper (PROTOCOL_silhouette_extension.md: the number of groups could not be estimated reliably); the selection task is what GUSTO-I applies.
% Cut 2026-10-14 (D-20261014-1 S03, F): CLAIMS C-48
% Cut 2026-10-15 (D-20261015-1, needs Tak's veto): 'and by the probability that the k retained regions include at least one discordant region' (second performance measure, reported only in the deleted fixed-n paragraph) -> dropped, AUC abbreviation defined here; 'Giving every measure the same k removes threshold choice from the comparison, though supplying k is itself an oracle' -> one clause; two performance sentences joined; 'the reported target is the middle of the three examined' kept
The measures are scored against the construction label on a selection task. An anchor and nine candidate regions are drawn; each measure ranks the candidates by their distance to the anchor and retains the $k$ closest, where $k$ is the number of true matches; supplying $k$ removes threshold choice from the comparison but is an oracle that the application does not use (Section~\ref{sec:application}). Performance is summarized by the area under the receiver operating characteristic curve (AUC) for discriminating true matches from discordant candidates (chance = 0.5), over 10{,}000 replications per cell at the base grid ($n = 25$--100 per region) and 3{,}000 at the extended grid ($n = 150$--2{,}000). The extended grid, the accuracy target and the definitions of the two behaviours below were fixed after the base-grid results had been seen, and the Monte Carlo intervals of Table~\ref{tab:study2_required_n_mc} were added later; the reported target is the middle of the three examined.

\subsubsection{Two Behaviours: Blind and Recovering}\label{sec:study1_classification}

% Paragraph reviewed 2026-09-27 (Option C, Harvey's recommendation under Tak's standing instruction; v6 paragraph review, queue #15): KS evidence given as a selection AUC (was a clustering ARI); RV1's above-chance plateau in the symmetric-severity cell explained as sampling-variance noise rather than resolved structure; operational labels kept as computed by R/required_n.R.
% Paragraph revised 2026-10-11 (D-20261011-11, cold referee M3/M5; Tak's standing instruction of 2026-10-11 23:49, individual veto open until submission): blindness grounded in the population argument for all seven cells
% Cut 2026-10-12 (D-20261012-3 length plan S03): 'Adding the SD as a further summary ...' sentence (no reported output behind it)
% Cut 2026-10-15 (D-20261015-1, needs Tak's veto): the enumeration of the SMD-blind cells ('counting each world's combined cell: every cell of Set 3, the two mean-matched cells of Set 4, the scale cell of Set 1, and the shape cell of Set 2') -> 'marked in the table' (the table marks the blind cells); count, AUC range and the population-SMD argument kept
Extending the sample-size grid to $n = 2{,}000$ separates the measures into two behaviours (Figure~\ref{fig:study2_auc}). A measure is blind in a cell when its AUC stays at chance through $n = 2{,}000$: the population value of its distance carries no information about the construction label, so the task is unsolvable for that measure at any sample size, an identification failure and not a power failure. This is exact by construction in every mean-matched cell, where the population SMD is zero. On the selection task, SMD is blind in \numval{nCellsSMDblind} of the \numval{nCellsSelection} cells (marked in Table~\ref{tab:study2_required_n_mc}; AUC \numval{aucBlindSMDmin}--\numval{aucBlindSMDmax} across the full grid in all of them).

% Paragraph revised 2026-10-11 (D-20261011-11, cold referee Minor 1; Tak's standing instruction of 2026-10-11 23:49, individual veto open until submission): weakest Set 4 cell is asymmetric severity
% Cut 2026-10-12 (D-20261012-3 length plan S21): 'Its failures in what follows are failures of efficiency' sentence
% Cut 2026-10-15 (D-20261015-1, needs Tak's veto): '(aucKSSet4symsevN2000 in the symmetric-severity cell)' -> dropped (the value of the second-weakest cell); 'identification succeeds, and what remains is a required-sample-size question' kept in shorter form
A measure is recovering when its AUC reaches the target and stays above it at every larger sample size: identification succeeds, and what remains is a required-sample-size question. $W_1$ and the KS statistic are recovering in every cell; the KS statistic is blind nowhere, reaching selection AUC \numval{aucKSSet4asymsevN2000} at $n = 2{,}000$ even in its weakest cell, asymmetric severity in Set 4.

\begin{figure}[ht]
\centering
\includegraphics[width=\textwidth]{fig_required_n_auc.pdf}
\caption{Selection AUC (true matches versus discordant candidates) as a function of per-region sample size ($n = 25$--2{,}000, log scale) for three distance measures in three cells of Study 1: the Set 3 combined cell (mean- and SD-matched mixtures), the Set 4 symmetric-severity cell (mean-matched displaced extremes), and the Set 4 bulk-shift control cell (uniform location shift of the whole population). The dashed horizontal line marks AUC = 0.9; chance corresponds to 0.5.}
\label{fig:study2_auc}
\end{figure}

\subsubsection{Selection: Required Sample Sizes}\label{sec:study1_selection}% Cut 2026-10-15 (D-20261015-1, needs Tak's veto): heading 'Selection: Required Sample Sizes and Error Rates' -> 'Selection: Required Sample Sizes' (the fixed-n error-rate paragraph is deleted)

% Paragraph reviewed 2026-09-27 (Option B, Harvey's recommendation under Tak's standing instruction; v6 paragraph review, queue #16): "decisive cells" replaced by all cells of Sets 3 and 4 (two omitted cells added to the table); the W1 column described as W1's requirement, not task difficulty.
% Paragraph revised 2026-10-11 (D-20261011-11, cold referee M3; Tak's standing instruction of 2026-10-11 23:49, individual veto open until submission): all 13 cells reported for W1 vs KS; headline 11 of 13 at AUC 0.9; target dependence disclosed
% Paragraph revised 2026-10-12 (D-20261012-2, Louis C4/M2/M10/m6/A9d; individual veto open until submission): Table 1 folded into Table 2 (13 cells x W1/KS/SMD); the "uniform shift suits the CDF supremum" mechanism removed (Sets 1-2 refute it) and replaced by the facts; "three by construction" attached to 11 of 13; the AUC 0.8 order stated with its overlapping intervals; SMD where not blind reported; 25-fold generated
% Cut 2026-10-14 (D-20261014-1 S04, R2c): CLAIMS C-08
% Cut 2026-10-14 (D-20261014-1 S18, R2c): CLAIMS C-49
% Cut 2026-10-15 (D-20261015-1, needs Tak's veto): 'and W1's own requirement ranges across the task from reqnW1Set4symsev (symmetric severity) to reqnW1maxSets34b (shape-symmetric mixtures) in Sets 3 and 4', the worked example of the symmetric-severity cell (reqnW1Set4symsevB, reqnKSSet4symsev, reqnRatioKSoverW1Set4symsev), the location-cell values (reqnW1Set1loc ... reqnKSSet2loc) and 'where neither measure is favoured by construction' -> dropped (values are in the table); the 'place ... place ... and overlap' chain rewritten as one sentence; the Set 4 anchor check (null SD ratios, paired AUC difference) condensed to the facts of the registered check (insurance, Louis): no mechanism, no 'the control is reported as a loss for W1'
Table~\ref{tab:study2_required_n_mc} reports, for all \numval{nCellsSelection} cells, the per-region sample size at which each measure's selection AUC reaches 0.9, with 95\% Monte Carlo intervals. At the AUC 0.9 target the 95\% intervals place $W_1$ below the KS statistic in \numval{nCellsW1FewerKS} cells, three of them the Set 4 subgroup cells, where the advantage follows from the construction (Section~\ref{sec:study1_design}), and the others including all six cells of Sets 1 and 2, which contain no outlying subgroup; they place the KS statistic below $W_1$ in \numval{nCellsKSFewer}, the bulk-shift control, and overlap in \numval{nCellsOverlap}, the shape-symmetric mixture (\numval{reqnW1Set3shapesym} [\numval{reqnW1Set3shapesymLo}, \numval{reqnW1Set3shapesymHi}] against \numval{reqnKSSet3shapesym} [\numval{reqnKSSet3shapesymLo}, \numval{reqnKSSet3shapesymHi}]). In the bulk-shift control the KS statistic reaches the target first (\numval{reqnKSSet4bulk} versus \numval{reqnW1Set4bulk}); in a check registered after an exploratory look (protocol, with the prior look disclosed, and outputs in the repository), for the same two-unit shift at $n = 200$ per region, the Set 4 anchor population multiplied the null sampling standard deviation of $\widehat{W}_1$ by \numval{set4NullSdRatioW1n200} (Monte Carlo SE \numval{set4NullSdRatioW1n200Mcse}) relative to $N(50, 10^2)$ and left that of the distribution-free KS statistic unchanged (\numval{set4NullSdRatioKSn200}, SE \numval{set4NullSdRatioKSn200Mcse}), and the paired AUC difference between $W_1$ and the KS statistic changed sign, from \numval{set4DeltaAucGaussN200d2} (SE \numval{set4DeltaAucGaussN200d2Mcse}) under the Gaussian anchor to \numval{set4DeltaAucMixN200d2} (SE \numval{set4DeltaAucMixN200d2Mcse}) under the Set 4 anchor.
% Sentences revised 2026-10-12 (D-20261012-3, REF M2; Louis four-eyes round 3): the mechanism of the control reversal stated from the registered check (results/set4_control_mechanism.csv, results/set4_null_sd_ratios.csv) within what PROTOCOL_set4_control_mechanism.md licenses — no proportionality to the integral, no generalization beyond the tested populations; 'rather than reversing the order in general' removed.
% Cut 2026-10-15 (D-20261015-1, needs Tak's veto): the sentence on SMD where it is not blind ('it reaches the target with sample sizes close to W1's in the Gaussian and log-normal location cells and with more elsewhere'; the SMD column of the table carries all 13 cells) and the values reqnSMDSet1comb, reqnW1Set1comb, reqnSMDSet2comb, reqnW1Set2comb -> dropped; 'the KS statistic discriminating better at the smallest sample sizes and W1 at the largest' and the SMD location-cell values (reqnSMDSet1loc, reqnW1Set1loc, reqnSMDSet2loc, reqnW1Set2loc) dropped; the crossing and the target dependence (AUC 0.95 / 0.8) kept
In the shape-symmetric mixture the two AUC curves cross near the target, so the overlap marks the crossing rather than equal performance: $W_1$ reaches AUC 0.95 first with separated intervals, while at AUC 0.8 the point estimates favour the KS statistic and the intervals overlap.

% Table 1 (tab:study2_required_n, Sets 3-4, W1/KS/SMD) folded into this table 2026-10-12 (D-20261012-2, Louis M10): same W1/KS values, SMD column added for all 13 cells; one float fewer.
% Paragraph reviewed 2026-09-27 (Option B, Harvey's recommendation under Tak's standing instruction; v6 paragraph review, queue #17): caption states how plateau and required n are computed.
\begin{table}[ht]
\centering
% Cut 2026-10-15 (D-20261015-1, needs Tak's veto): caption of the required-n table: 'using the Monte Carlo standard errors of the two AUC estimates that bracket the crossing' dropped (pure deletion; the notes to the table state it) -> deleted (no replacement wording)
\caption{Per-region sample size at which selection AUC reaches 0.9, by cell and measure, interpolated on $\log n$ over the grid $n = 25$--2{,}000, with 95\% Monte Carlo intervals from the delta method applied to the log-$n$ interpolation. Cells whose AUC stays at chance (\numval{aucBlindSMDmin}--\numval{aucBlindSMDmax}) across the full grid are marked ``blind''; $\leq 25$ marks a target reached at the smallest grid size, where no interval is available, and such a cell counts as separated when 25 lies below the other measure's lower limit.\label{tab:study2_required_n_mc}}
\input{generated/tab_study2_required_n_mc}% GENERATED by R/make_paper_tables.R from results/required_n_selection.csv (complete tabular)
\begin{tablenotes}
% ERRFIX-2026-10-11 (Louis review B1): "combined" is a discordant role (C1/C2 in build_set3: asymmetric and bimodal at once), not a pool of the other cells
\item Cells are contrast types within each world; the ``combined'' cell holds discordant regions that differ from the anchor in both of the world's contrast types at once (in Set 3, asymmetric and bimodal), not a pool of the other cells.
\item Intervals reflect only the Monte Carlo error of the two AUC estimates that bracket the crossing, not the error of interpolating between grid points; limits are rounded outward to whole patients. ``Separate'' means the $W_1$ and KS intervals do not overlap. All measures are computed on the same replications, so non-overlap is a conservative indication of a difference and overlap does not show equal performance.
\end{tablenotes}
\end{table}

% Cut 2026-10-15 (D-20261015-1, needs Tak's veto): fixed-sample-size error-rate paragraph (n = 100 in Set 4: probability that the three retained regions include at least one discordant region, for W1 and KS, with chance level; 'three of the four discordant types favour W1 by construction') deleted (pure deletion; the definition of that second performance measure is dropped from Sec. 3.1.1) -> deleted (no replacement wording)

% Paragraph reviewed 2026-09-27 (Option B, Harvey's recommendation under Tak's standing instruction; v6 paragraph review, queue #19): all six discordant regions reported (was four); "clinically mildest/worst" replaced by position on the W1 scale; "admits" replaced.
% Paragraph deleted 2026-10-12 (D-20261012-2, Louis length cut / m11, P3): the Set 4 population-ordering paragraph (true W1 and KS of the six discordant regions; "no KS threshold retains exactly the regions with true W1 <= c") said of itself that it was not independent evidence; no abstract or Discussion claim rested on it (CLAIMS C-30 withdrawn). The values remain in results/selection_sim_summary.csv.

% Section 3.1.4 (clustering) removed 2026-09-27 (Tak): the number of groups could not be estimated reliably (PROTOCOL_silhouette_extension.md, criterion not met), so clustering is not a method that identifies pooling regions; not reported (process check, like the supplements).

% DECIDED (Tak, 2026-09-05): the outline's remaining Study 2 material will NOT
% be written -- (i) the rho = W1/(KS*sigma_EM) regime heuristic (per-cell
% rho_trans never entered the verification gate) and (ii) the oracle-best-
% threshold steelman with the Part 1B operating characteristics. The section
% is complete as it stands.
% DECIDED (Tak, 2026-08-30, option A): outline subsection 3.1.4 (procedure
% justification, anchor-vs-clustering, 100k reps) will NOT be written here --
% its content stays in the Discussion paragraph sourced from
% results/anchor_vs_clustering_tau_100k.csv, which is its home.

\subsection{Study 2: Estimation Properties}\label{sec:study2}

% Sentence cut 2026-10-12 (D-20261012-2, Louis length cut, P3): the "underwrite the calibration framework" sentence repeated the opening of Section 3.
% Cut 2026-10-12 (D-20261012-3 length plan S08): units/coverage-unitless sentence of the Study 2 opening (repeated in every caption)
% Cut 2026-10-14 (D-20261014-1 S01, R2d): CLAIMS none

\subsubsection{Design}\label{sec:sim_design}

% Paragraph reviewed 2026-09-27 (Option A, Harvey's recommendation under Tak's standing instruction; v6 paragraph review, queue #22): S6 described as scale combined with skewness; unsupported "most realistic" and "typical of ALT" claims removed; caption describes the table.
We designed seven systematic scenarios for methodological validation, examining controlled distributional differences in location, scale, scale combined with skewness (S6, a mean-matched log-normal), and a location--scale combination. Table~\ref{tab:scenarios} summarizes these scenarios with their true $W_1$ values, in the original units of the effect modifier.

\begin{table}[ht]
\centering
\caption{Study 2 simulation scenarios with true $W_1$ values.\label{tab:scenarios}}
\begin{tabular}{lllll}
\toprule
\textbf{ID} & \textbf{Description} & \textbf{Distribution 1} & \textbf{Distribution 2} & \textbf{True $W_1$} \\
\midrule
S1 & Null (identical) & $N(50, 10^2)$ & $N(50, 10^2)$ & \numval{trueW1S1} \\
S2 & Location 0.2$\sigma$ & $N(50, 10^2)$ & $N(52, 10^2)$ & \numval{trueW1S2} \\
S3 & Location 0.5$\sigma$ & $N(50, 10^2)$ & $N(55, 10^2)$ & \numval{trueW1S3} \\
S4 & Location 1.0$\sigma$ & $N(50, 10^2)$ & $N(60, 10^2)$ & \numval{trueW1S4} \\
S5 & Scale 1.5$\times$ & $N(50, 10^2)$ & $N(50, 15^2)$ & \numval{trueW1S5} \\
S6 & Skew + scale (Log-normal) & $N(50, 10^2)$ & LogN$(\numval{muLogS6tab}, 0.5^2)$ & \numval{trueW1S6} \\
S7 & Location + Scale & $N(50, 10^2)$ & $N(55, 15^2)$ & \numval{trueW1S7} \\
\bottomrule
\end{tabular}
\begin{tablenotes}
% Cut 2026-10-15 (D-20261015-1, needs Tak's veto): note to the scenario table: '(against 10 for Distribution 1, so the contrast is in scale as well as shape)', 'the table shows two decimals, and', S6/S7 parameter details and 'the share of intervals lying above the truth ... was computed with Monte Carlo truths for S6 and S7' dropped; the exact true W1 values of S5-S7 (trueW1S5exact, trueW1S6exact, trueW1S7exact) kept -> deleted (no replacement wording)
\item S6: $\sigma_{\log} = 0.5$ and $\mu_{\log} = \log 50 - \sigma_{\log}^2/2 = \numval{muLogS6note}$, giving mean 50, SD $\numval{sdS6}$, CV $= \numval{cvS6pct}\%$ and skewness $\numval{skewS6}$. True $W_1$ is exact for S1--S5 ($\sqrt{2/\pi}\,|\sigma_1 - \sigma_2|$ in S5) and obtained by numerical quadrature for S6 and S7; the performance measures use the exact values $\numval{trueW1S5exact}$ (S5), $\numval{trueW1S6exact}$ (S6) and $\numval{trueW1S7exact}$ (S7).
\end{tablenotes}
\end{table}

% Cut 2026-10-12 (D-20261012-3 length plan S18): R-version parenthetical (recorded in sessionInfo)
% Cut 2026-10-14 (D-20261014-1 S19, F): CLAIMS none
% Cut 2026-10-14 (D-20261014-1 S20, R2c): CLAIMS none (reproducibility)
% Cut 2026-10-15 (D-20261015-1, needs Tak's veto): 'Study 2 was run in R with the bootstrap in C++ via Rcpp, replications distributed over a PSOCK cluster with L'Ecuyer--CMRG random-number streams seeded per cell (clusterSetRNGStream, seed 42 + 1,000k ...); because replicates are assigned to workers dynamically, individual replicates are not bit-reproducible' -> moved to the Data availability statement (one sentence); 'Every replication produced a finite estimate and interval.' dropped; one clause pointing to the note of the estimation table (replaces the deleted four-item list of performance measures)
For each scenario, we generated samples of size $n = 50$, 100, and 200 per region. We performed 10{,}000 replications per scenario--sample size combination; the Monte Carlo standard error is at most \numval{mcseBoundProp} for any reported proportion, and over the 21 cells it is at most \numval{mcseBiasMax} for bias, \numval{mcseRmseMax} for RMSE and \numval{mcseWidthMax} for mean CI width. Bootstrap confidence intervals were computed as in Section~\ref{sec:estimation} with $B = 2{,}000$ resamples; the performance measures are defined in the note to Table~\ref{tab:estimation}.

% Cut 2026-10-15 (D-20261015-1, needs Tak's veto): the four-item list of operating characteristics (bias, RMSE, coverage with 'Under S1 (W_1 = 0) coverage is structurally zero', CI width) deleted (pure deletion; the note to the estimation table defines all four and the zero coverage under S1 is stated in Sec. 3.2.2) -> deleted (no replacement wording)

\subsubsection{Results}\label{sec:sim_results}

% Paragraph revised 2026-10-11 (D-20261011-11, cold referee Minor 3; Tak's standing instruction of 2026-10-11 23:49, individual veto open until submission): bias decay described as observed (n^-1/2), not as slow
% Cut 2026-10-14 (D-20261014-1 S08, R2c): CLAIMS C-17, C-23, C-32
% Cut 2026-10-15 (D-20261015-1, needs Tak's veto): '(scenarios are specified with sigma = 10)' and 'and the bias fell at a rate consistent with n^{-1/2}: at n = 200 it is about half that at n = 50' dropped (the rate statements are removed from the whole of Sec. 3.2.2); the next paragraph (S3-S7 bias at n = 200 and S4 within Monte Carlo error; S1-S2 larger relative to the true value) joined to this one -> deleted (no replacement wording)
Table~\ref{tab:estimation} presents bias, coverage, RMSE and mean CI width of $\widehat{W}_1$ across scenarios and sample sizes, in the original units of the effect modifier. The estimator showed positive bias under the null scenario (S1), decreasing from \numval{biasS1n50} at $n=50$ to \numval{biasS1n200} at $n=200$: the finite-sample bias of Section~\ref{sec:estimation} ($E\,\widehat{W}_1 \geq W_1$), at its largest when the true value equals zero ($F_1 = F_2$ almost everywhere). The estimator remains consistent ($\widehat{W}_1 \to 0$ as $n \to \infty$ under the null). For the non-null scenarios away from the boundary (S3--S7), bias was small relative to each scenario's true $W_1$ and diminished with sample size: at $n = 200$ it was at most \numval{biasMaxS3toS7n200} (S7), and in S4 it was within Monte Carlo error at every $n$; near the boundary (S1, S2) it was larger relative to the true value.

\begin{table}[ht]
\centering
\caption{Bias, coverage of the 95\% percentile bootstrap confidence interval, RMSE and mean CI width of $\widehat{W}_1$ by scenario and sample size, in the original units of the effect modifier.\label{tab:estimation}}
\input{generated/tab_estimation}%  GENERATED by R/make_paper_tables.R from results/w1_raw_summary.csv (complete tabular)
\begin{tablenotes}
\item Bias is $\text{Mean}(\widehat{W}_1) - W_1$ and RMSE is $\sqrt{\text{Mean}[(\widehat{W}_1 - W_1)^2]}$; coverage is the proportion of the 95\% percentile bootstrap intervals that contain $W_1$, and is zero under S1 because $W_1 = 0$ lies at the parameter space boundary and a percentile interval cannot contain it; CI width is the mean of (upper $-$ lower). Scenarios are specified with $\sigma = 10$. Results based on 10{,}000 replications with $B = 2{,}000$ bootstrap resamples.
\end{tablenotes}
\end{table}

% Paragraph revised 2026-10-11 (D-20261011-11, cold referee Minor 3; Tak's standing instruction of 2026-10-11 23:49, individual veto open until submission): relative, not slow
% Paragraph revised 2026-10-12 (D-20261012-2, Louis C3 / Mike A1, m7; individual veto open until submission): the n^{-1/2} rate holds only under the null (S1); off the boundary the bias falls faster (bias(50)/bias(200): S2 5.7, S5 4.4, S6 3.4, S7 3.3, against 2 for n^{-1/2}; results/derived/study2_bias_ratio.csv) and is within Monte Carlo error from n = 100 in S3 and S4; the S4 "confirming" sentence cut
% Cut 2026-10-15 (D-20261015-1, needs Tak's veto): paragraph on bias away from the boundary ('Under the null (S1) ... the bias fell at the n^{-1/2} rate; in the other non-null scenarios ... it fell faster, the bias at n = 200 being at most one third of that at n = 50') joined to the preceding paragraph, rate statements dropped (both); 'Figure 2(A)' -> the estimation table
% Sentence corrected 2026-10-12 11:00 (Katrina four-eyes, results/VERIFY_FOUREYES_20261012.md): S3's bias at n = 100 is 3.7 MCSE, so "within Monte Carlo error" holds for S4 only; "close to n^{-1}" replaced by the observed range of ratios (3.3-5.7 against 4).

% Cut 2026-10-15 (D-20261015-1, needs Tak's veto): Figure 2 (fig:simulation; panels A-C: bias, coverage, mean CI width, the same quantities as the estimation table) removed; references to it repointed to the estimation table (tab:estimation) -> deleted (no replacement wording)

% Paragraph revised 2026-10-11 (D-20261011-11, cold referee M2; Tak's standing instruction of 2026-10-11 23:49, individual veto open until submission): near-nominal replaced by the range and the MCSE count; coverage tied to the standardized distance
% Paragraph revised 2026-10-12 (D-20261012-2, Louis C2 / Mike A8; individual veto open until submission): the standardized-distance sentence was contradicted by results/w1_raw_summary.csv (n = 100: S3 3.66 SD -> 0.951, S5 3.84 -> 0.917, S7 4.15 -> 0.935); coverage follows bias relative to the sampling SD (results/derived/study2_bias_over_sd.csv); S5 range added
% Cut 2026-10-14 (D-20261014-1 S10, R2c): CLAIMS C-09, C-17
% Cut 2026-10-15 (D-20261015-1, needs Tak's veto): the bias-over-SD annotations (biasOverSdS5n100, biasOverSdS7n100, biasOverSdS3n100, biasOverSdS6n100), the reference to the intervals table and the S4 coverage-range sentence (covS4rangeAllN, 'with no degradation at larger n, within two Monte Carlo standard errors of 0.95 only at n = 200') dropped; the coverage range with the cell count, the S5/S7 upward-displacement statement, the S6 / S4-100 two-sided statement and the S5/S2 n = 50 values kept; the two coverage paragraphs joined -> deleted (no replacement wording)
For the non-null scenarios S3--S7 at $n \geq 100$ (Table~\ref{tab:estimation}), coverage ranged over \numval{covRangeS3toS7N100up2dp}; of these ten cells, \numval{nCellsCovBelow2MCSE} lie more than two Monte Carlo standard errors below 0.95, the lowest being S5 at $n = 100$ (\numval{covS5n100}). The shortfall is largest in S5 and S7, whose CDFs cross; there the intervals lie above the true value, in line with the upward bias of Section~\ref{sec:estimation}. In S6 and in S4 at $n = 100$, whose bias relative to the sampling standard deviation of $\widehat{W}_1$ is smaller or within Monte Carlo error, the shortfall is small and the misses fall on both sides, so a given $n$ does not by itself fix the coverage. S5 (scale) and S2 (small location shift, near the boundary) showed undercoverage at $n=50$ (\numval{covS5n50} and \numval{covS2n50}) that improved to \numval{covS5n200} and \numval{covS2n200} at $n=200$.

% Miss sides: R/study1_miss_side.R -> results/study1_miss_side.csv (2026-09-27).
% Paragraph reviewed 2026-09-27 (Option B, Harvey's recommendation under Tak's standing instruction; v6 paragraph review, queue #23): bias sentence (inconsistent with consistency) removed; bootstrap inconsistency at F1 = F2 cited; miss side of the undercovering intervals reported; verdict wording removed.
% Paragraph revised 2026-10-11 (D-20261011-11, cold referee Minor 5; Tak's standing instruction of 2026-10-11 23:49, individual veto open until submission): conservativeness scoped to small distances
% Cut 2026-10-14 (D-20261014-1 S12, R2c): CLAIMS C-17
% Cut 2026-10-14 (D-20261014-1 S13, R2c): CLAIMS C-17
% Cut 2026-10-14 (D-20261014-1 S22, R2d): CLAIMS C-22, C-23
% Cut 2026-10-15 (D-20261015-1, needs Tak's veto): '(Table 5, Figure 2(B))' pointers and 'This does not indicate estimator failure---the estimator is consistent---but' reordered; 'the interval misses from below ... against 2.5% nominal' and the S6 exception kept; 'conservative reading' scope unchanged -> deleted (no replacement wording)
Coverage for S1 (null) is zero because its true $W_1$ lies at the parameter space boundary: a percentile interval cannot contain a true value of zero at any sample size, because with continuous data every bootstrap replicate of $\widehat{W}_1$ is positive; whether an observed distance is distinguishable from no difference is read against the null floor. Where coverage falls short in the non-null scenarios (S2, S5 and S7 at smaller $n$), the interval misses from below: the true value exceeded the upper limit in at most \numval{missAboveMaxS2S5S7pct}\% of replications in S2, S5 and S7 at any sample size (never in S2), against 2.5\% nominal. The upper confidence limit is therefore a conservative reading of how large a difference could be in S2, S5 and S7; this is not a property of every scenario whose CDFs cross, since in S6 the true value exceeded the upper limit at about the nominal rate.
% Sentence revised 2026-10-12 (D-20261012-2, Louis C2): "away from the boundary it is nominal rather than conservative" removed — S7 (0.58 SD) still misses from below (results/study1_miss_side.csv).

% Cut 2026-10-15 (D-20261015-1, needs Tak's veto): 'A difference between two cells is within simulation noise when it is smaller than about zDiffFactor times the standard error' dropped; the post hoc basic (reflected) interval analysis (its formula, its coverage range covBasicRangeS3toS7N100up, its S1 and S2 values covBasicS1n50 and covBasicS2n50, 'this bias correction trades coverage away from the boundary for coverage at it') -> one clause (insurance, Louis); references to the intervals table dropped
Monte Carlo standard errors for the reported coverage probabilities are at most \numval{mcseBoundPropB} at 10{,}000 replications (largest observed \numval{mcseCovMax}), and a single cell is consistent with nominal coverage when it lies within about twice its own standard error of 0.95. Because the application reads an unresolved distance at its upper limit, the one-sided coverage of the 97.5\% upper limit, $P(W_1 \leq U)$, is the operating characteristic that matters there: for S3--S7 at $n \geq 100$ it ranged over \numval{covUpperMinS3toS7N100up}--\numval{covUpperMaxS3toS7N100up} against a nominal 0.975: above nominal where the interval lies above the truth (S5, S7, and S3 at $n = 100$), near nominal in S6 and in S3 at $n = 200$, and below it in S4 at $n = 100$. A bias-corrected (basic) interval, examined post hoc, covered \numval{covBasicRangeS3toS7N100up} in S3--S7 from $n = 100$ and exceeded the percentile interval only under S1 and in S2 at $n = 50$ (\numval{covBasicS2n50} against \numval{covS2n50}; repository).
% Sentences added 2026-10-12 (D-20261012-3, EXT M4(c)(d)): one-sided upper-limit coverage and the basic interval from results/study2_interval_variants.csv (post hoc, labelled).

% Cut 2026-10-15 (D-20261015-1, needs Tak's veto): the table of coverage of the percentile interval, of its one-sided 97.5% upper limit and of the basic interval (tab:study2_intervals) removed; the percentile and upper-limit coverage are in the text and the estimation table, the basic-interval range in the repository -> deleted (no replacement wording)

% Paragraph and table added 2026-10-13 (D-20261012-3, REF2 M5 / EXT M8): registered extension of Study 2 (PROTOCOL_study2_extension.md, designed after the external review and before any run; results/w1_raw_ext_summary.csv)
% Cut 2026-10-14 (D-20261014-1 H5, R2c): CLAIMS C-09; Harvey menu 5
% Cut 2026-10-15 (D-20261015-1, needs Tak's veto): the registration history in parentheses ('its first eight cells were registered after the external review and before the run, and three further cells ... as an amendment registered before those cells ran'; kept in the note to the extension table), '(the cells are described in the note to Table 7)' and 'lowest for the shift on the t_3 reference (covExtE2b, with a positive bias of biasExtE2b)' (in the table) -> dropped from the text; the coverage range, MCSE, number of replications and the two operating-point values kept
A registered extension of Study 2 (protocol and outputs in the repository) added cells outside the original design: unequal regional sizes, a heavy-tailed reference distribution, integer-grid data, $n = 1{,}000$ per region and the application's operating point (Table~\ref{tab:study2_ext}). Across the ten non-null cells coverage ranged over \numval{covExtMin}--\numval{covExtMax} (Monte Carlo standard error at most \numval{mcseCovExtMax}; \numval{nRepsExtMin}--10{,}000 replications per cell); at the application's operating point it was \numval{covExtE5} with equal sizes and \numval{covExtE6} with unequal, rounded samples, with upper-limit coverage \numval{covUpperExtE5} and \numval{covUpperExtE6}.

\begin{table}[ht]
\centering
\caption{Registered extension of Study 2: unequal regional sizes, a $t_3$ reference distribution, integer-grid data, $n = 1{,}000$ per region and the application's operating point (0.2 SD). Bias, RMSE and mean width in the original units; coverage of the 95\% percentile interval and of its one-sided 97.5\% upper limit.\label{tab:study2_ext}}
\input{generated/tab_study2_ext}%  GENERATED by R/make_paper_tables.R from results/w1_raw_ext_summary.csv (complete tabular)
\begin{tablenotes}
\item Cells: E1, unequal regional sizes (100 against 300; 200 against 1{,}000); E2, a $t_3$ reference scaled to SD 10, under the null and under a 5-unit shift; E3, integer-grid data (the S3 contrast rounded to integers at $n = 100$, 200 and 1{,}000); E4, $n = 1{,}000$ per region for the S3 and S5 contrasts; E5 and E6, the application's operating point (the 0.2 SD contrast of S2 at $n = 1{,}000$ per region, and at 2{,}900 against 1{,}200 with both samples rounded to multiples of 10, as systolic blood pressure is recorded).
\item Under the $t_3$ null the true $W_1$ is zero, so percentile coverage is structurally zero and upper-limit coverage identically 1, as under S1. The rounded-data truth is the $W_1$ of the rounded population distributions. Replications per cell as shown; Monte Carlo standard error of coverage at most \numval{mcseCovExtMax}. Protocol \texttt{PROTOCOL\_study2\_extension.md}, registered after the external review; the two operating-point cells and the integer-grid cell at $n = 1{,}000$ were added as an amendment before they were run.
\end{tablenotes}
\end{table}

% Paragraph reviewed 2026-09-27 (Option C, Harvey's recommendation under Tak's standing instruction; v6 paragraph review, queue #24): units corrected (S3 is read on the age scale only for the example); "proportionally" corrected; the single Delta_max example kept here (duplicate in the S3 paragraph deleted) and tied to what n >= 100 does not guarantee.
% Cut 2026-10-12 (D-20261012-3 length plan S14): arithmetic of the Delta_max width example and the propagation sentence (ratio to the margin kept)
% Cut 2026-10-12 (D-20261012-3 length plan S14): arithmetic of the Delta_max width example and the propagation sentence (ratio to the margin kept)
% Cut 2026-10-14 (D-20261014-1 S15, R2c): CLAIMS C-32
% Cut 2026-10-14 (D-20261014-1 S16, R2c): CLAIMS C-32
% Cut 2026-10-15 (D-20261015-1, needs Tak's veto): 'at the O(n^{-1/2}) rate expected from Wasserstein plug-in theory', the S3 RMSE values (rmseS3n50, rmseS3n200), 'Mean CI width narrowed with n (Figure 2(C))' and the CI-width range (ciwRangeNonNullN100) dropped; the Delta_max width example with its ratio to the margin unchanged -> deleted (no replacement wording)
RMSE decreased with increasing $n$ (Table~\ref{tab:estimation}). For S3 on the age scale with the illustrative bound $L_{\text{UB,age}} = 0.01$/yr of the GUSTO-I application, a $\widehat{W}_1$ CI width of \numval{ciwS3n100} years at $n = 100$ corresponds to a $\Delta_{\max}$ interval about \numval{dmaxWidthS3n100overMargin} times the 1\%pt margin used in Section~\ref{sec:app_clinical}, so $n = 100$ per region does not by itself make the $\Delta_{\max}$ interval narrow relative to a clinical margin.

% S1 95th percentiles: results/study1_miss_side.csv (q95_est), from results/w1_raw_simulation.rds.
% Paragraph reviewed 2026-09-27 (Option B, Harvey's recommendation under Tak's standing instruction; v6 paragraph review, queue #25): "null hypothesis" -> null scenario; repeated means removed; closing aligned with the approved Section 2.3 reading (at or below, not evidence of a difference) and scoped to sigma = 10, equal n.
% Cut 2026-10-12 (D-20261012-3 length plan S15): S1 mean-CI-width sentences (in the precision table)
Under the null scenario (S1), the informative summary is the upper tail rather than the mean: the 95th percentile of $\widehat{W}_1$ under S1 is \numval{q95S1n50}, \numval{q95S1n100} and \numval{q95S1n200} at $n = 50$, 100 and 200 for this scenario ($\sigma = 10$, equal $n$), and an observed distance at or below the corresponding value is not evidence of a difference.

% NOTE (2026-08-30): the former SMD sensitivity table (tab:smd, S3/S5/S6
% detection contrast) was deleted per the restructure outline -- Study 2's
% six-method decision comparison supersedes it. Discussion paragraphs still
% citing S5/S6 as SMD-blindness evidence carry queued review edits (decision 4).
% Paragraph revised 2026-10-11 (D-20261011-11, cold referee Minor 4; Tak's standing instruction of 2026-10-11 23:49, individual veto open until submission): relative bias scoped to S3--S7
% Summary list deleted 2026-10-12 (D-20261012-2, Louis length cut, P3): the three-item "In summary" list repeated the bias, coverage and precision paragraphs number for number; the bridge to Delta_max below carries the summary. RMSE at n = 200 ranged from rmseMinN200 (S5) to rmseMaxN200 (S1); both keys remain in the registry.

% Paragraph deleted 2026-09-27 (Option B, Harvey's recommendation under Tak's standing instruction; v6 paragraph review, queue #26): duplicated the Delta_max example of the CI-width paragraph (P3).

% Paragraph reviewed 2026-09-27 (Option B, Harvey's recommendation under Tak's standing instruction; v6 paragraph review, queue #27): "reliable estimation and inference" narrowed to bias and coverage; fixed-L qualifier added; section pointer corrected.
% Paragraph revised 2026-10-11 (D-20261011-11, cold referee M2; Tak's standing instruction of 2026-10-11 23:49, individual veto open until submission): near-nominal replaced by the range with its scale
% Paragraph revised 2026-10-12 (D-20261012-2, Louis M1/M9/m2; individual veto open until submission): the MCSE count carried here; "estimation criterion" replaced by "estimation statement"; "about 0.4 SD or more"
% Cut 2026-10-12 (D-20261012-3 length plan S16): 'In summary' paragraph (coverage range restated a fourth time)


%==============================================================================
\section{Application: Hypothetical Thrombolytic MRCT Using GUSTO-I Data}\label{sec:application}
%==============================================================================

% Paragraph reviewed 2026-09-27 (Option A, Harvey's recommendation under Tak's standing instruction; v6 paragraph review, queue #28): "per-EM" spelled out; stand-in assumption for the regional populations stated.
% Cut 2026-10-14 (D-20261014-1 A01, F): CLAIMS C-10, C-39
% Cut 2026-10-15 (D-20261015-1, needs Tak's veto): 'a hypothetical scenario grounded in publicly available data from GUSTO-I' kept; 'A sponsor developing a novel thrombolytic agent (Drug T) for AMI plans a Phase 3 MRCT across multiple regions and uses IPD ...' -> 'A sponsor planning a Phase 3 MRCT of a novel thrombolytic agent (Drug T) uses IPD ...'; 'to quantify how the distributions of candidate effect modifiers in other regions differ from those in a designated anchor region' kept
We illustrate the per-effect-modifier Wasserstein-1 framework through a hypothetical scenario grounded in publicly available data from GUSTO-I, a large international trial in acute myocardial infarction (AMI). A sponsor planning a Phase~3 MRCT of a novel thrombolytic agent (Drug~T) uses IPD from GUSTO-I ($N = \numval{gustoN}$; 16 regions)\cite{gusto1993} to quantify how the distributions of candidate effect modifiers in other regions differ from those in a designated anchor region, treating the GUSTO-I regions as stand-ins for the regional populations the new trial would enrol; GUSTO-I was not designed as an MRCT and its regions are anonymized.

\subsection{Scenario Design and Effect Modifier Selection}\label{sec:app_scenario}

% Paragraph reviewed 2026-09-27 (Option B, Harvey's recommendation under Tak's standing instruction; v6 paragraph review, queue #29): age and SBP given the same class-level evidence statement (FTT: benefit irrespective of age and blood pressure); repeated "Drug T is a novel agent" removed.
% Cut 2026-10-14 (D-20261014-1 A02, R2c): CLAIMS C-25
% Cut 2026-10-15 (D-20261015-1, needs Tak's veto): the second 'chosen for illustration', 'Class-level indirect evidence can at most inform a sponsor-specified upper bound (Section 2.4)' and 'so class transferability would not yield a defensible numerical prior for Drug T' -> dropped; 'The evidence available for such an agent would not permit numerical estimation of L_clinical' and the following 'We therefore treat L_clinical as unknown a priori ... apply the L* reverse-calculation pathway to both' joined into one paragraph; the FTT finding (condition restored 2026-10-12) unchanged
For this illustrative exercise, the sponsor takes age and systolic blood pressure (SBP) as the candidate effect modifiers for Drug~T. The Fibrinolytic Therapy Trialists' (FTT) collaborative meta-analysis of existing fibrinolytic agents\cite{ftt1994} concluded that, among patients presenting with ST elevation or bundle-branch block, benefit is present irrespective of age and of blood pressure, and reports no per-unit CATE slope on the absolute mortality scale for either, so the evidence available would not permit numerical estimation of $L_{\text{clinical}}$; we therefore treat $L_{\text{clinical}}$ as unknown a priori for both candidate effect modifiers and apply the $L^*$ reverse-calculation pathway (equation~\ref{eq:lstar}) to both.

% Paragraph reviewed 2026-09-27 (Option A, Harvey's recommendation under Tak's standing instruction; v6 paragraph review, queue #30): no change to the text; the pair-level IQR column removed from the Region 8 table (orphaned by the queue #34 deletion).
% Paragraph revised 2026-10-11 (D-20261011-11, cold referee Minor 12; Tak's standing instruction of 2026-10-11 23:49, individual veto open until submission): data screening stated
% Cut 2026-10-14 (D-20261014-1 A03, R2c): CLAIMS C-10
% Cut 2026-10-14 (D-20261014-1 A10, F): CLAIMS C-01
% Cut 2026-10-15 (D-20261015-1, needs Tak's veto): data-screening paragraph (w1MaxChangeExclAge, w1MaxChangeExclSbp, w1MaxChangeWinsorSbp, 'these checks were made after the main analysis', the share of readings on a multiple of 10 mmHg in the 16 regions, 'the step CDFs of Appendix A.1 apply') -> insurance D (Louis): records retained, effect of excluding them on null-floor classification and SBP ranks, winsorizing (post hoc), heaping most in Region 8; 'representing for instance a regulatory market where local efficacy evidence is required', 'and the anchor role, not its size, is what the exercise illustrates', 'GUSTO-I regions are larger than many in a MRCT; Study 2 examines smaller regional sizes' and 'a target region compared with all other regions pooled is defined in the same way, with the target excluded from the reference' dropped; the pointer to the Region 8 baseline table (removed) dropped
The sponsor selects Region~8 as the anchor region ($n = \numval{nR8}$ patients). Records with SBP $= 0$~mmHg (\numval{nSbpZero} patients) or age $> 100$ years (\numval{nAgeOver100} patients) were retained; excluding them changed no null-floor classification below and \numval{nRankChangesExcl} adjacent SBP ranks, and winsorizing at the 1st and 99th percentiles changed \numval{nRankChangesWinsor} of the 30 ranks; both checks were made post hoc. SBP is heaped on multiples of 10~mmHg, most in Region~8 (\numval{sbpDigit0R8}\% of readings), so recording practice may contribute to the SBP distances (not quantified here).
% Sentences added 2026-10-12 (D-20261012-3, EXT M7(c) / REF Minor 9): results/gusto_sensitivity_outliers_summary.csv, results/gusto_sbp_digit_preference.csv (post hoc sensitivity, labelled in the provenance).

% Cut 2026-10-15 (D-20261015-1, needs Tak's veto): table of Region 8 baseline characteristics (tab:gusto_r8) removed; the anchor's n stays in the text -> deleted (no replacement wording)

\subsection{Null Floor: Which Distances Are Distinguishable from No Difference?}\label{sec:app_operability}

% Paragraph reviewed 2026-09-27 (Option B, Harvey's recommendation under Tak's standing instruction; v6 paragraph review, queue #31): "unresolved" attached to a distance, not a partner; R8 introduced; forward pointer to the distances; unused null-mean column removed.
% Paragraph revised 2026-10-11 (D-20261011-11, cold referee M4; Tak's standing instruction of 2026-10-11 23:49, individual veto open until submission): null floor named as a test
% Cut 2026-10-14 (D-20261014-1 A04, R2c): CLAIMS C-12, C-42
% Cut 2026-10-14 (D-20261014-1 A11, R2d): CLAIMS C-12, C-42
% Cut 2026-10-14 (D-20261014-1 A12, R2c): CLAIMS C-12
% Cut 2026-10-14 (D-20261014-1 H2b, R2c): CLAIMS C-12
% Cut 2026-10-15 (D-20261015-1, needs Tak's veto): 'Table 7 summarises the result' (tab:operability, removed) -> the two sentences of the following paragraph (six of fifteen age distances unresolved at level 0.05; fourteen of fifteen SBP distances resolved, R2 the exception) moved here; 'Before interpreting the distances of Section 4.3', 'applied separately to each pair', 'which is exact under exchangeability', '(2,000 permutations)', 'so the anchor-resampled floor is retained as the display reference', the Bonferroni count nPermBonf and 'drawing the null from the partner's own distribution instead of the anchor's leaves every distance on the same side of its 95th percentile' dropped; the permutation agreement count kept
We computed the null floor of Section~\ref{sec:estimation} separately for each candidate partner, from two independent resamples, with replacement, of Region~8's (R8) empirical distribution at sizes $n_{\text{R8}} = \numval{nR8floor}$ and $n_{\text{partner}}$, recomputing $\widehat{W}_1$ over $2{,}000$ repetitions. We call an observed distance unresolved when it lies at or below the 95th percentile of its null distribution (a one-sided test of $F_1 = F_2$ at level 0.05). A permutation test of $F_1 = F_2$ on the pooled sample classifies all 30 distances in the same way (\numval{nPermAgree} of 30 agree). On age, the distances of six of the fifteen partners---R1, R4, R5, R7, R9 and R15---are unresolved at level 0.05: they are ones the estimator would produce with appreciable probability from identical distributions. On systolic blood pressure, by contrast, fourteen of fifteen are resolved, the single exception being R2.
% Sentence added 2026-10-12 (D-20261012-3, EXT M5 / REF Minor 4): results/gusto_null_floor_permutation_summary.csv.

% Cut 2026-10-15 (D-20261015-1, needs Tak's veto): table of the null floor across partners, by effect modifier (tab:operability) removed; its content (six of fifteen age distances, one of fifteen SBP distances unresolved) is in the preceding paragraph -> deleted (no replacement wording)

% Cut 2026-10-15 (D-20261015-1, needs Tak's veto): 'The two candidate effect modifiers differ in how many distances clear the null floor, and the difference matters for how the subsequent distances should be read.' (opening) -> dropped; the two counts it introduced are in the preceding paragraph

% Cut 2026-10-14 (D-20261014-1 A05, F): CLAIMS C-12
% Cut 2026-10-14 (D-20261014-1 H2a, F): CLAIMS C-12; Harvey menu 2

% Paragraph reviewed 2026-09-27 (Option B, Harvey's recommendation under Tak's standing instruction; v6 paragraph review, queue #33): defensive opening removed; "plug-in estimator" wording; the upper-limit reading tied to the Study 1 miss-side result.
% Paragraph revised 2026-10-11 (D-20261011-11, cold referee M2; Tak's standing instruction of 2026-10-11 23:49, individual veto open until submission): the n >= 100 transfer removed; null-floor resolution and coverage share the standardized distance
% Paragraph revised 2026-10-12 (D-20261012-2, Louis C2 inheritance / length cut; individual veto open until submission): "governed by the same quantity" removed (the standardized distance does not govern coverage); defensive opening sentence cut
% Cut 2026-10-14 (D-20261014-1 A06, R2c): CLAIMS C-12
% Cut 2026-10-15 (D-20261015-1, needs Tak's veto): 'The small age distances of the six partners unresolved at level 0.05 are therefore not evidence of a difference;' -> 'How large the differences of the six partners unresolved at level 0.05 could be is read from ...'; 'Neither resolution against the null floor nor the coverage of the percentile interval is fixed by the sample size alone; ... the partner samples range from partnerNmin to partnerNmax observations, but for a distance at or below the null floor the interval is read at its upper limit whatever the sample size.', 'so the intervals of Table 8 are read with that evidence rather than with the original cells' and 'Study 2 scenarios (Section 3.2.2)' pointer reworded -> dropped; the sentence tying the GUSTO-I distances (gustoDistSdMin-gustoDistSdMax SD) to the Study 2 operating-point cells kept
How large the differences of the six partners unresolved at level 0.05 could be is read from the upper confidence limits in Table~\ref{tab:gusto_nabcd}, which for small distances near the boundary err above rather than below the true value in the Study~2 scenarios (Section~\ref{sec:sim_results}). The GUSTO-I distances are \numval{gustoDistSdMin}--\numval{gustoDistSdMax} SD of Region~8's distributions, below the \numval{minW1SdS3toS7}~SD at which S3--S7 begin and in places above the 0.2~SD of S2; the extension cells at this operating point (Table~\ref{tab:study2_ext}) gave percentile coverage \numval{covExtE5} and \numval{covExtE6} and upper-limit coverage \numval{covUpperExtE5} and \numval{covUpperExtE6}, and the intervals of Table~\ref{tab:gusto_nabcd} are read with that evidence.
% Sentence added 2026-10-13 (D-20261012-3, REF2 M5 / Louis round 3 M12): results/derived/gusto_dist_sd.csv, results/w1_raw_ext_summary.csv.

\subsection{Distributional Assessment: Per-Effect-Modifier $W_1$ Across 15 Partner Regions}\label{sec:app_nabcd}

% Reopened 2026-09-27 (Tak): the simulation identifies the method, GUSTO-I applies it; partners ordered by W1 (ranks in the table).
% Paragraph revised 2026-10-11 (D-20261011-11, cold referee M1/M3; Tak's standing instruction of 2026-10-11 23:49, individual veto open until submission): headline 11 of 13; attainment named
% Paragraph revised 2026-10-12 (D-20261012-2, Louis C1/M2/M3 / Mike A3; individual veto open until submission): "at the GUSTO-I region sizes both suffice" removed (KS asymmetric severity 1,252 [1,226, 1,278] vs R16 at 1,231; required n is specific to the simulated contrasts); "three by construction"; attainment within the class
% Cut 2026-10-14 (D-20261014-1 A07, R2c): CLAIMS C-11
% Cut 2026-10-15 (D-20261015-1, needs Tak's veto): 'Study 1 showed that W1 and the KS statistic both reach selection AUC 0.9 in every simulated world (Section 3.1.3);' (restates Sec. 3.1.2), '; W1 is used because its heterogeneity bound has a per-unit clinical slope as its constant (Section 2.2)', 'in the original measurement units (years for age, mmHg for SBP)' (in the caption of the table), 'We therefore computed' and 'Figure 3 (a forest plot)' dropped -> deleted (no replacement wording)
Study~1's required sample sizes refer to its simulated contrasts and do not transfer to the GUSTO-I distances, which are read against their null floors (Section~\ref{sec:app_operability}). We computed $\widehat{W}_1$ between Region~8 and each of the remaining 15 regions for both candidate effect modifiers, ordered the partners by it, and give percentile bootstrap 95\% confidence intervals ($B = 2{,}000$) in Table~\ref{tab:gusto_nabcd}.

\begin{table}[ht]
\centering
\caption{Per-effect-modifier Wasserstein-1 distance $\widehat{W}_1$ in original variable units (years for age, mmHg for SBP) with 95\% percentile bootstrap CIs for Region~8 versus each partner region (GUSTO-I), listed by region number, with each partner's rank by $\widehat{W}_1$ within each candidate effect modifier.\label{tab:gusto_nabcd}}
\input{generated/tab_gusto_nabcd}%  GENERATED by R/make_paper_tables.R from results/gusto_r8_w1_per_pair.csv (complete tabular)
\begin{tablenotes}
\item Point estimates with 95\% percentile bootstrap CIs ($B = 2{,}000$, seed = 2026). Regions are listed by region number; Rank orders the partners by $\widehat{W}_1$ within each candidate effect modifier (1 = closest to Region~8); the bracketed rank interval is the 2.5th--97.5th percentile of the rank over 1{,}000 patient-level bootstrap replicates that resample within every region (Region~8 shared across partners; seed 20261011), and adjacent ranks whose intervals overlap are not resolved against each other. $^\dagger$Distances unresolved against their null floor (Section~\ref{sec:app_operability}) are not ordered among themselves and share the point-rank range shown; the bracketed interval that follows it is that partner's own bootstrap rank interval. R6 and R5, adjacent in the SBP ranking, differ in $\widehat{W}_{1,\text{SBP}}$ by \numval{w1SbpR6minusR5}~mmHg, so their ranks are effectively tied.% Cut 2026-10-15 (D-20261015-1, needs Tak's veto): table reference 'Table 7' (tab:operability, removed) -> 'Section 4.2' (same content in the text); 'which for these partners reaches the ninth or tenth position' dropped
% Note revised 2026-10-12 (D-20261012-2, Louis m12): the orphan R6/R5 sentence given its point; the two CI-overlap sentences, superseded by the rank intervals, cut.
\end{tablenotes}
\end{table}

% Cut 2026-10-15 (D-20261015-1, needs Tak's veto): forest plot of the W1 point estimates and intervals (fig:gusto_forest; the same numbers as the table of per-partner distances) removed; the dagger convention is stated in the note to that table -> deleted (no replacement wording)

% Paragraph reviewed 2026-09-27 (Option B, Harvey's recommendation under Tak's standing instruction; v6 paragraph review, queue #34): cross-modifier IQR comparison deleted (contradicts the approved Section 2.2 scope statement); "similarity" -> difference.
% Cut 2026-10-14 (D-20261014-1 A14, F): CLAIMS C-11
% Cut 2026-10-14 (D-20261014-1 H3, F): CLAIMS C-11; Harvey menu 3

% Paragraph revised 2026-10-11 (D-20261011-11, cold referee M4; Tak's standing instruction of 2026-10-11 23:49, individual veto open until submission): R16 example replaced; ranks read with their resampling intervals
% Cut 2026-10-14 (D-20261014-1 A08, F): CLAIMS C-43
% Cut 2026-10-15 (D-20261015-1, needs Tak's veto): the two paragraphs on rank uncertainty and on the R2 / R9 contrast joined into one; 'where uncertainty intervals overlap substantially', 'and both of its rank intervals span most of the middle of the table', and the forward pointer 'Section 5 reports the rank agreement across the 15 partners' dropped; R6's ranks, the most distant partners (R2 on age, R9 on SBP), the four distances of the reversal and 'the reversal rests on the two large ones' kept -> deleted (no replacement wording)
Several partners fall in the middle of each ranking, where a rank is not resolved against its neighbours: the rank intervals of Table~\ref{tab:gusto_nabcd}, obtained by resampling patients within every region, show that R6, ranked ninth on age and fifth on SBP, can move across most of the middle of the table, whereas the most distant partners, R2 on age and R9 on SBP, have rank intervals confined to the most distant positions, so a bare rank is not read as a resolved order. R2 has the second-largest $\widehat{W}_{1,\text{age}}$ (\numval{w1AgeR2} years) but the smallest $\widehat{W}_{1,\text{SBP}}$ (\numval{w1SbpR2b} mmHg), while R9 is among the six smallest $\widehat{W}_{1,\text{age}}$ (\numval{w1AgeR9} years) but the largest $\widehat{W}_{1,\text{SBP}}$ (\numval{w1SbpR9b} mmHg). Both small distances are unresolved against their null floors (Section~\ref{sec:app_operability}), so the reversal rests on the two large ones and illustrates that the distance on one candidate effect modifier need not carry information about the distance on another.

\subsection{Clinical Interpretation}\label{sec:app_clinical}

% Paragraph reviewed and approved by Tak 2026-09-06 (Option C of the paragraph
% review; simulated-review point R1.2(a), application side): aligned with the
% Section 2.4 judgment framing; the "supports a conservative upper bound ...
% exceeding the empirical prognostic gradients" phrasing retired; notation
% unified to L_{UB,em} across Section 4.4 in the same pass.
% Paragraph reviewed 2026-09-27 (Option B, Harvey's recommendation under Tak's standing instruction; v6 paragraph review, queue #36): "We require" removed (the L* pathway does not require a bound); duplication of Section 2.4 cut; "approximately" removed (exactly 10%pt per decade).
% Cut 2026-10-14 (D-20261014-1, needs Tak's veto): 'Following Section 2.4, we read L* against illustrative bounds that are the sponsor's clinical inputs, informed by the class-level evidence cited there: L_UB,age = 1e-2 per year (10%pt per decade) and L_UB,SBP = 2e-3 per mmHg (2%pt per 10 mmHg) on the 30-day mortality scale.' -> paragraph deleted (the same two bounds are stated in Sec. 2.4 and in the note of Table 9; the references to 'Section 4.4' for the bounds are repointed to Section 2.4)


% Paragraph reviewed 2026-09-27 (Option C, Harvey's recommendation under Tak's standing instruction; v6 paragraph review, queue #37): interval definition moved to Section 2.4 (queue #8), reference kept; "plausible" -> illustrative; lower-interval reading made explicit; link to the null floor added.
% Cut 2026-10-15 (D-20261015-1, needs Tak's veto): '; the factor would equal 1 only if the partner's CDF never crossed the anchor's' (parenthetical) and 'Each interval is read against the illustrative bound ...' as a separate sentence dropped (joined to the first); the three readings of the L* interval shortened and the monotone shrink-factor ranges kept -> deleted (no replacement wording)
Table~\ref{tab:gusto_lstar_joint} reports, for each partner region, $L^*$ with its percentile interval (Section~\ref{sec:inference}), read against the illustrative bound $L_{\text{UB}}$ for its candidate effect modifier. Where the whole interval lies above $L_{\text{UB}}$, even the upper confidence limit of the distance would require a slope above the bound to produce $\Delta_{\text{clin}} = 1$\%pt; where it contains $L_{\text{UB}}$, the data are compatible with a difference that could and one that could not produce it; where it lies below $L_{\text{UB}}$, a slope within the bound could produce it even at the lower confidence limit of the distance. Only the lower end of the interval, which comes from the upper confidence limit of $\widehat{W}_1$, carries the reading for distances unresolved against the null floor. If the CATE were additionally assumed monotone, equation~(\ref{eq:mono_bound}) would replace $\Delta_{\max}$ by the fraction $(W_1 + |\mu_1 - \mu_2|)/(2 W_1)$ of it, which ranges from \numval{monoRatioAgeMin} to \numval{monoRatioAgeMax} across the age distances and from \numval{monoRatioSbpMin} to \numval{monoRatioSbpMax} across the SBP distances (point estimates, including the unresolved R2 on SBP).
% Sentence added 2026-10-12 (D-20261012-3, REF M1): the monotone-bound factor in the application (results/derived/gusto_monotone_bound_ratio.csv).

% Paragraph reviewed 2026-09-27 (Option B, Harvey's recommendation under Tak's standing instruction; v6 paragraph review, queue #38): caption describes the table (bounds moved to a note); unresolved distances flagged; R1 SBP lower limit disambiguated; brackets unified with Table 5.
% Table rows: R/gusto_lstar_table.R -> results/gusto_lstar_intervals.csv (from results/gusto_r8_w1_per_pair.csv), 2026-09-26.
\begin{table}[ht]
\centering
\caption{$L^*$ at $\Delta_{\text{clin}} = 1$\%pt with 95\% percentile interval for all 15 partner regions, listed by region number. $L^*$ is the minimum clinical slope (per year for age, per mmHg for SBP) required for the observed distributional difference to produce a treatment effect difference of $1$\%pt on the absolute 30-day mortality scale.\label{tab:gusto_lstar_joint}}
\input{generated/tab_gusto_lstar_joint}%  GENERATED by R/make_paper_tables.R from results/gusto_lstar_intervals.csv + results/gusto_operability.csv (complete tabular)
\begin{tablenotes}
\item Regions are listed by region number; their ranks by $\widehat{W}_1$ are given in Table~\ref{tab:gusto_nabcd}. $^\dagger$Distance unresolved against its null floor (Section~\ref{sec:app_operability}). The illustrative bounds of Section~\ref{sec:inference} are $L_{\text{UB,age}} = 1\times10^{-2}$ /yr and $L_{\text{UB,SBP}} = 2\times10^{-3}$ /mmHg; the lower limit of R1's SBP interval, \numval{lstarSbpLowerR1}, lies just above $L_{\text{UB,SBP}}$.% Cut 2026-10-15 (D-20261015-1, needs Tak's veto): table reference 'Table 7' (tab:operability, removed) -> 'Section 4.2'
\end{tablenotes}
\end{table}

%==============================================================================
\section{Discussion}\label{sec:discussion}
%==============================================================================

% Paragraph reviewed and approved by Tak 2026-08-30 (Option B of the paragraph
% review; resolves queued edit item 4 for this paragraph -- decision 4, 2026-08-15).
% Paragraph reviewed 2026-09-27 (Option C, Harvey's recommendation under Tak's standing instruction; v6 paragraph review, queue #39): "adds an inferential result" -> characterizes its estimation; constant described as sponsor-specified; (iii) scoped to the Study 1 scenarios; Komiyama citation moved to the RV clause.
% Paragraph revised 2026-10-11 (D-20261011-11, cold referee M1/M2/Minor 9; Tak's standing instruction of 2026-10-11 23:49, individual veto open until submission): attainment in the summary; sample size attached to recovery; coverage range with its scale
% Paragraph revised 2026-10-12 (D-20261012-2, Louis M8/m17/M1/m2; individual veto open until submission): the enumeration aligned with the three requirements of Section 2.1 (coverage is the estimation characterization, not requirement (iii)); the 95-word sentence split; the MCSE count carried; "about 0.4 SD or more"
% Cut 2026-10-12 (D-20261012-3 length plan D13): KS total-variation clause (Methods 2.1-2.2)
% Cut 2026-10-14 (D-20261014-1, needs Tak's veto): Its estimation is characterized by Study 2: percentile bootstrap intervals whose coverage was [range] for true distances of about [x] to [y] SD in the location, scale, and combined scale-and-skewness scenarios at n >= 100 per region, below 0.95 beyond Monte Carlo error in [k] of the ten cells, with a characterized boundary behavior near the null (Section 2.3) -> Its estimation is characterized by Study 2 (Section 3.2.2); the coverage range and the boundary behaviour near the null are reported there
% Cut 2026-10-15 (D-20261015-1, needs Tak's veto): the Study 1 recovery range ('at n = reqnW1minSets34--reqnW1maxSets34 per region in Sets 3 and 4'), the coverage range covRangeS3toS7N100up2dp (the count nCellsCovBelow2MCSE and the SD scope of the true distances kept), 'with a characterized boundary behaviour near the null', 'an identification failure that no sample size repairs' and the section pointers -> dropped; 'SMD supports no bound valid across the Lipschitz class, since regions with equal means can have W_1 > 0 and hence, by equation (dual), regional treatment effects differing by up to L W_1' dropped; 'informed by clinical evidence', 'a treatment effect change per unit of the effect modifier' and 'and characterizes the estimation of its answer' dropped
This paper answers the requirements set out in Section~\ref{sec:existing}. The $W_1$ distance meets requirement (i): in Study 1, SMD is at chance at every sample size examined in \numval{nCellsSMDblind} of the \numval{nCellsSelection} selection cells, whereas $W_1$ reaches the discrimination target in all four worlds. It meets requirements (ii) and (iii): the heterogeneity bound \(\Delta_{\max}\) (equation~\ref{eq:delta_max}) translates distributional distance into a clinically interpretable scale, is sharp when the slope is the only restriction on the CATE and tightens under monotonicity (equation~\ref{eq:mono_bound}), and its constant is a clinical slope that the sponsor specifies. Its estimation is characterized by Study~2: for true distances of about \numval{minW1SdS3toS7} to \numval{maxW1SdS3toS7} SD at $n \geq 100$ per region, percentile bootstrap coverage was below 0.95 beyond Monte Carlo error in \numval{nCellsCovBelow2MCSE} of the ten cells (Section~\ref{sec:sim_results}).
% Restored 2026-10-14 01:30 (Louis review-fix of the cut, K3): the D01 cut removed the MCSE count and the SD scope (2026-10-12 corrections, rule 2).

% Paragraph reviewed 2026-09-27 (Option B, Harvey's recommendation under Tak's standing instruction; v6 paragraph review, queue #40): procedural "must be evaluated" and "require" replaced by what the distances show; binary "similar" wording removed.
% Cut 2026-10-12 (D-20261012-3 length plan D01): restatement of L* monotonicity/comparison (Methods 2.4) compressed
% Cut 2026-10-12 (D-20261012-3 length plan D15): worked-example gloss (Sec. 4 opening)
% Cut 2026-10-14 (D-20261014-1 D02, R2c): CLAIMS C-35, C-06
% Cut 2026-10-15 (D-20261015-1, needs Tak's veto): 'illustrated by the GUSTO-I example (Section 4)' kept; 'so the distance on a modifier left out of the assessment need not be inferable from those included' and the lead clause of 'Second, W1 ranking alone does not convey clinical weight: ...' and 'and both are reported' dropped (the R2 / R9 contrast itself is in Sec. 4.3) -> deleted (no replacement wording)
These findings carry two implications for pooling assessments, illustrated by the GUSTO-I example (Section~\ref{sec:application}). First, when several candidate effect modifiers are considered, the per-modifier distances need not agree: a partner close to the anchor on one effect modifier can be distant on another, as in the contrast between R2 and R9. Second, within one effect modifier $L^* = \Delta_{\text{clin}} / W_1$ is a strictly decreasing function of $W_1$, so calibration does not reorder partners; it compares each partner's $L^*$ with the bound $L_{\text{UB}}$ specified for that modifier.

% Paragraph reviewed 2026-09-27 (Option B, Harvey's recommendation under Tak's standing instruction; v6 paragraph review, queue #41): partner-count recommendation ("no more than four") removed per the quantify-not-decide principle; duplicate "120 region pairs" removed; correlation stated as weak (r now reported in Supplement C).
% Reopened 2026-09-27 (Tak): Komiyama et al. introduce one example approach; no "recipe", no comparison.
% Paragraph revised 2026-10-11 (D-20261011-11, cold referee M4/M5; Tak's standing instruction of 2026-10-11 23:49, individual veto open until submission): patient-level bootstrap interval for the rank agreement; clustering omission referenced
% Cut 2026-10-12 (D-20261012-3 length plan D02): clustering pointer shortened (guard-pinned clause 'not reported' kept)
% Cut 2026-10-12 (D-20261012-3 length plan D03): R2/R9 restatement (para 2 and Sec. 4.3)
% Cut 2026-10-12 (D-20261012-3 length plan D04): 'interval reflects sampling uncertainty within this trial' clause
% Cut 2026-10-14 (D-20261014-1 D03, R2c): CLAIMS C-11
% Cut 2026-10-15 (D-20261015-1, needs Tak's veto): the Komiyama / cluster-analysis restatement ('These implications also bear on the cluster analysis of regions that Komiyama et al. introduce ... the clustering task run under our protocol is not reported') deleted (pure deletion; the clustering task is declared in Sec. 3.1.1); 'How many pooled regions to form, and how orderings that differ across modifiers are reconciled, remain clinical judgments.' dropped (pure deletion; the paper's closing paragraph states that it informs, not prescribes); the Spearman correlation of age and SBP distances, 'no positive agreement' and 'aggregated over the two modifiers could mask' kept -> deleted (no replacement wording)
Across the 15 partners of Table~\ref{tab:gusto_nabcd}, the orderings by age and by systolic blood pressure distance show no positive agreement (Spearman correlation $\numval{spearmanAgeVsSbpSigned}$; 95\% interval \numval{spearmanAgeVsSbpBootLo} to \numval{spearmanAgeVsSbpBootHi} from the patient-level bootstrap of Table~\ref{tab:gusto_nabcd}; it is computed on the point ordering of all 15 partners, including the distances left unordered by the null floors); the interval excludes the strong positive agreement under which one modifier's distances could stand in for the other's, so a single distance aggregated over the two modifiers could mask differences that a per-modifier assessment retains.

% Paragraph reviewed and approved by Tak 2026-08-30 (Option B of the paragraph
% review; resolves queued edit item 5 -- decision 4, 2026-08-15). This paragraph
% is the Discussion's sole carrier of the GUSTO agreement finding, per that
% decision (the abstract stays silent on it).
% Paragraph reviewed 2026-09-27 (Option B, Harvey's recommendation under Tak's standing instruction; v6 paragraph review, queue #42): "three gaps" aligned with Section 2.1 (limitations of the existing measures); closing claim rested on the Study 2 efficiency evidence instead of the KR-duality tautology.
% Paragraph revised 2026-10-11 (D-20261011-11, cold referee M1/M3/M5; Tak's standing instruction of 2026-10-11 23:49, individual veto open until submission): constant described; survey cited; headline 11 of 13 with attainment and construction caveat
% Paragraph revised 2026-10-12 (D-20261012-2, Louis M5/M6/m10 / length cut; individual veto open until submission): the restatement of Section 2.1 compressed to one sentence; patient-level bootstrap intervals for the W1-vs-SMD Spearman (results/gusto_rank_spearman_boot.csv); IST-1 and IST-3 cited; survey correlations to three decimals; "like the KS statistic" attached to recovery
% Cut 2026-10-14 (D-20261014-1, needs Tak's veto): Of the limitations set out in Section 2.1, the location-only limitation of SMD is structural, and Study 1 exhibits its consequence: in the mean-matched cells, SMD is at chance at every sample size examined while W1 recovers the partner structure (Section 3.1.2); the lack of a clinical scale is addressed by ...; and the limitations of KL divergence are avoided by construction ... Addressing these limitations does not, however, mean returning a different answer in every dataset -> one sentence with the same three limitations (the Study 1 restatement of paragraph 1 (i) dropped)
% Cut 2026-10-15 (D-20261015-1, needs Tak's veto): the first sentence ('Of the limitations set out in Section 2.1, the location-only limitation of SMD is structural, the lack of a clinical scale is addressed by the heterogeneity bound, and those of KL divergence are avoided by construction ...'), the survey of region pairs in GUSTO-I, IST-1 and IST-3 (surveyPairs, surveySpearmanMin-Max, surveyRows; W1-KS agreement surveySpearmanW1KSmin-Max), the efficiency tally against the KS statistic ('fewer patients than the KS statistic in nCellsW1FewerKS of nCellsSelection cells ... heavy-tailed ...'), 'and its contribution there was the clinical calibration rather than a different answer', 'this does not mean returning a different answer in every dataset' and 'The reason to use W1 is therefore not efficiency against the KS statistic' deleted (the survey is declared in the Data availability statement, the tally is in Sec. 3.1.3 with its construction caveat); kept: the GUSTO-I Spearman correlations of W1 with the SMD with their intervals, the SD ratio, 'configurations beyond location cannot be excluded before the distances are computed', the value of W1 (not blind where the SMD is; clinical-scale translation) -> deleted (no replacement wording)
In GUSTO-I, whose regional differences are predominantly differences in location (partner standard deviations within a factor of \numval{sdRatioMax} of the anchor's), $W_1$ ordered the 15 partners much as the SMD did (Spearman correlation \numval{spearmanW1SMDage}, 95\% interval \numval{spearmanW1SMDageLo} to \numval{spearmanW1SMDageHi}, for age and \numval{spearmanW1SMDsbp}, \numval{spearmanW1SMDsbpLo} to \numval{spearmanW1SMDsbpHi}, % Cut 2026-10-15 (D-20261015-1, Louis review-fix, needs Tak's veto): survey result direction disclosed (selective-reporting guard; protocol 6)
for systolic blood pressure, from the patient-level bootstrap of Table~\ref{tab:gusto_nabcd}), and a registered survey of region pairs in GUSTO-I, IST-1 and IST-3 found agreement of the same order (Spearman correlation \numval{surveySpearmanMin}--\numval{surveySpearmanMax}; repository).\cite{ist1_2000,ist3_2012} The limitations of the SMD bind where distributional differences go beyond location---mixtures and rare displaced subgroups---and configurations beyond location cannot be excluded before the distances are computed. The reason to use $W_1$ is that, like the KS statistic, it is not blind where the SMD is, and that it carries a translation to the clinical scale whose constant is a per-unit clinical slope, valid for every CATE in the class and conservative for any particular one.
% Sentences revised 2026-10-12 (D-20261012-3, REF M2 / REF M1): efficiency tally no longer the reason; survey W1-KS agreement (results/real_data_survey_summary.csv, spearman_w1_ks) reported; attainment wording replaced.

% Paragraph reviewed and approved by Tak 2026-08-31 (Option B of the paragraph
% review; resolves queued edit item 6 -- found by Louis during the item-5 review).
% The evidence for the identification advantage lives in the preceding
% paragraph and Section 3; this paragraph deliberately names no moments.
% Notation-only edit 2026-09-27 (decision round): bare L -> L_clinical; "CATE sensitivity evidence" -> evidence on the clinical slope.
% Paragraph revised 2026-10-12 (D-20261012-2, Louis length cut, P3): the restatement of the two calibration pathways of Section 2.4 cut; the separability statement kept.
% Cut 2026-10-14 (D-20261014-1, needs Tak's veto): 'The distributional assessment and its calibration are separable: the identification and efficiency results of the preceding paragraph are properties of the distance itself and hold whether or not L_clinical is available; clinical calibration adds an outcome-scale reading but does not change them.' -> paragraph deleted (Harvey menu 4; Sec. 2.4 states the two calibration pathways and Sec. 5 paragraph 1 states what the distance does without them)


% Paragraph deleted 2026-09-27 (Option B, Harvey's recommendation under Tak's standing instruction; v6 paragraph review, queue #43): restated the procedure of Section 2.4 (P3) and carried "L is estimable" against the judgment framing.

% Paragraph reviewed 2026-09-27 (Option B, Harvey's recommendation under Tak's standing instruction; v6 paragraph review, queue #44): "distinctive advantage" removed (SMD needs the same inputs); representativeness assumption for external sources added; regulatory and Matsushima sentences removed as repeats of the Introduction.
% Cut 2026-10-15 (D-20261015-1, needs Tak's veto): 'The assessment can therefore be made at the planning stage, before any regional data from the new trial exist.' merged into the first sentence; 'like other covariate-distribution measures,', 'Computed from baseline data alone,' and the parenthetical '(region means alone give only |mu_1 - mu_2| <= W_1)' dropped; the last two sentences joined ('complements rather than replaces assessments based on observed treatment effects' for 'Computed from baseline data alone, it complements assessments based on observed treatment effects and does not replace them') -> deleted (no replacement wording)
Because $W_1$ requires only baseline effect modifier distributions, regional data can be assembled from prior trials, registries, electronic health records, or other real-world evidence (RWE) sources, provided the source population approximates the population each region is expected to enrol, so that the assessment can be made at the planning stage. It uses neither treatment assignment nor outcomes, requires individual-level baseline values, and complements rather than replaces assessments based on observed treatment effects.

% Paragraph reviewed 2026-09-27 (Option C, Harvey's recommendation under Tak's standing instruction; v6 paragraph review, queue #45): common-CATE assumption of Proposition 1 added as a limitation (extrinsic factors); estimator notation; S5 undercoverage and its miss side stated; n >= 100 given as the Study 1 estimation criterion.
% Paragraph revised 2026-10-11 (D-20261011-11, cold referee M2; Tak's standing instruction of 2026-10-11 23:49, individual veto open until submission): estimation criterion carries its scale
% Cut 2026-10-15 (D-20261015-1, needs Tak's veto): 'biological plausibility, clinical relevance, and benefit--risk balance extend beyond what any single distributional measure can quantify' -> 'it is one input within a holistic regulatory evaluation' only; 'The framework has further limitations.' and 'the framework operates on whatever candidates the sponsor specifies' dropped; the binary / ordinal parenthetical dropped (pure deletion)
The $W_1$ distance is one input within a holistic regulatory evaluation.\cite{iche17,matsushima2024} It treats continuous effect modifiers marginally, leaving categorical and multivariate extensions for future work, and it is stated for collapsible effect scales (Section~\ref{sec:methods}). Which baseline characteristics are effect modifiers is taken as given; their identification requires evidence on treatment-by-covariate interaction rather than on prognosis alone (Section~\ref{sec:inference}) and lies outside the framework.
% Restored 2026-10-14 01:30 (Louis review-fix of the cut, K4): the H1 and D08 cuts together removed the only statement of this scope.
% Cut 2026-10-15 (D-20261015-1, needs Tak's veto): the S2 and S5 n = 50 coverage values (covS2n50disc, covS5n50disc), 'estimates near the null are read against the null floor' and 'A systematic measurement difference of size b between regions changes W_1 by at most |b|' -> dropped, no remaining carrier in the paper (the n = 50 values are in Sec. 3.2.2; 'estimates near the null are read against the null floor' is in Sec. 2.3); 'such as standard of care or concomitant therapy' dropped; the registered-extension sentence, the common-CATE assumption and the unmeasured-modifier sentence kept in shorter form
Because $W_1$ is non-negative, $\widehat{W}_1$ exhibits positive bias near zero, zero percentile coverage when the true value is zero (S1), and undercoverage near the boundary and where the bias is large relative to the sampling standard deviation, with the intervals erring above the true value (Section~\ref{sec:sim_results}).\cite{andrews2000} The coverage reported for Study~2 from $n = 100$ per region refers to true distances of about \numval{minW1SdS3toS7} to \numval{maxW1SdS3toS7} SD of the reference distribution, with equal regional sizes; the registered extension (Table~\ref{tab:study2_ext}) covers 0.2 SD at the application's sample sizes, unequal sizes, a heavy-tailed reference and grid-recorded data, but not every combination the application presents. Proposition~\ref{prop:heterogeneity} assumes a CATE common to all regions, so regional differences in extrinsic factors that alter $\tau$ itself lie outside what $W_1$ measures, and unmeasured effect modifiers may contribute to heterogeneity beyond what the distance captures.

% Paragraph reviewed and approved by Tak 2026-09-22 (Option B of the paragraph
% review; response-queue item (1) paragraph D, the Discussion receiver promised by
% Section 2.2 for the threshold-CATE case and by the M1 rewrite for the KS
% alternative): L_UB as judgment, average slope as lower bound, threshold CATE
% as design (worst case over the class), KS bound as the alternative, locality
% of the Lipschitz requirement. The former one-sentence L limitation in the
% preceding paragraph was removed as duplicated here.
% Paragraph reviewed 2026-09-27 (Option B, Harvey's recommendation under Tak's standing instruction; v6 paragraph review, queue #46): "declaring regions dissimilar" replaced (the paper declares nothing); locality statement made precise, including its limit under full support; additive per-modifier reading stated.
% Cut 2026-10-14 (D-20261014-1, needs Tak's veto): First, ... L_UB is specified by the sponsor, ... conditional on it; reporting L* leaves the comparison open to any alternative bound. Second, an L_clinical taken from an estimated average interaction slope is a lower bound on the Lipschitz constant ...; the safety factor of Section 2.4 is a judgment of the same kind as margin discounting, and the framework does not supply its value. Third, ... -> First (merged with Second) ...; Second, ... [two sentences merged, 'Third' renumbered]
% Cut 2026-10-14 (D-20261014-1, needs Tak's veto): 'The requirement is also local: ... for distributions with full support, such as normal distributions, that range is the whole line. Finally, the bound is per effect modifier: it assumes the CATE varies with that modifier alone. When the CATE is additive across modifiers, ..., the per-modifier bounds combine as ..., and each per-modifier Delta_max is the contribution attributable to that modifier; under interaction between modifiers they do not, because ...' -> same locality statement and full-support qualifier ('although' replaces the sentence break), 'it assumes the CATE varies with that modifier alone' dropped (it stays in Proposition 1 and the abstract), the two additive sentences joined into one with all formulas and qualifiers kept (C-21 sole carrier)
% Cut 2026-10-15 (D-20261015-1, needs Tak's veto): 'The calibration rests on the Lipschitz premise of Proposition 1, and its limitations are those of the premise.' (lead sentence), 'the same holds for an L_clinical taken from an estimated average interaction slope, a lower bound on the Lipschitz constant, for which the safety factor of Section 2.4 is a judgment the framework does not supply', 'informed but not determined by prior evidence' and 'the Kolmogorov--Smirnov alternative for that case is stated in Section 2.2' (sentence J there) dropped; the locality statement (the bound holds with L replaced by the Lipschitz constant on {x : F_1 neq F_2}; full-support qualifier), the additive-CATE formula |tau_1 - tau_2| <= sum_k L_k W_1^(k) with 'each per-modifier Delta_max is the contribution attributable to that modifier', and the weighted-l1 joint bound |tau_1 - tau_2| <= W_d(P_1, P_2) 'with or without interaction' -> dropped; replaced by insurance I (Louis): per-modifier bounds sum to a lower bound on the joint one, W_d >= sum_k L_k W_1^(k), so under interaction the per-modifier assessment can understate heterogeneity; estimating W_d is outside this paper
The limitations of the calibration are those of the Lipschitz premise of Proposition~\ref{prop:heterogeneity}. First, the constant is a clinical judgment (Section~\ref{sec:inference}), so conclusions drawn from $\Delta_{\max}$, or from the comparison of $L^*$ with $L_{\text{UB}}$, are conditional on it. Second, for a CATE that changes steeply over a narrow range of the effect modifier---a threshold or step function---the Lipschitz constant is large and equation~(\ref{eq:heterogeneity_bound}) is correspondingly loose. This is by design: the bound is a worst case over the CATEs compatible with the stated $L_{\text{clinical}}$, and a loose bound overstates rather than understates the potential difference. % Cut 2026-10-15 (D-20261015-1, Louis review-fix, needs Tak's veto): W_d: premise and inequality restored ('the joint one' had no antecedent); additive case stated
Finally, the bound is per effect modifier. If the CATE is Lipschitz with respect to $d(x, y) = \sum_k L_k |x_k - y_k|$, then $|\bar{\tau}_1 - \bar{\tau}_2| \leq W_d(P_1, P_2)$, the Wasserstein distance between the joint distributions under cost $d$, and $W_d \geq \sum_k L_k W_1^{(k)}$; the per-modifier sum is a valid bound for a CATE additive across modifiers but, under interaction, can understate heterogeneity. Estimating $W_d$ is outside this paper.
% Sentences added 2026-10-12 (D-20261012-3, EXT M6): the weighted-l1 joint bound and its relation to the marginal sum (both inequalities verified by Louis and Mike).

% Paragraph reviewed 2026-09-27 (Option B, Harvey's recommendation under Tak's standing instruction; v6 paragraph review, queue #47): "where none was previously available" and "filling the gap left by ICH E17" narrowed (P5; Komiyama is an existing recipe; E17 is non-prescriptive by design); contribution statement split into a closing paragraph.
% Paragraph revised 2026-10-12 (D-20261012-2, Louis M7; individual veto open until submission): the outcome-Lasso "front-end" removed (it selects prognostic factors, which Section 2.4 distinguishes from effect modification; in Komiyama et al. the Lasso weights distances rather than selecting modifiers)
% Cut 2026-10-14 (D-20261014-1 D08, R2d): CLAIMS none
Inference tailored to the boundary behavior of $\widehat{W}_1$---resampling fewer than $n$ observations\cite{bickel1997} or the numerical delta method for directionally differentiable functionals\cite{hong2018}---could extend the operational range of the intervals at small distances.

% Paragraph revised 2026-10-11 (D-20261011-11, cold referee M1; Tak's standing instruction of 2026-10-11 23:49, individual veto open until submission): attainment in the closing sentence
% Cut 2026-10-12 (D-20261012-3 length plan D07): duplicate 'operationalizing the similarity' clause in closing sentence
The principal contribution of this paper is a distributional similarity measure for effect modifiers whose heterogeneity bound holds for every CATE with a given per-unit clinical slope, is sharp over that class and conservative for any particular CATE, and whose calibration constant is that slope. With bootstrap confidence intervals and clinical calibration via $\Delta_{\max}$ or $L^*$, $W_1$ supplies a quantitative input to the judgment of ``similar enough'', informing, not prescribing, pooling decisions.

%==============================================================================
% Back matter
%==============================================================================

\bmsection*{Author contributions}

[Author 1] conceived the research question and led the project. [Author 2] developed the mathematical framework and conducted the simulation study. [Author 3] contributed to the application example and manuscript preparation. All authors reviewed and approved the final manuscript.

\bmsection*{Acknowledgments}

The authors thank [colleagues] for helpful discussions on ICH E17 implementation.

\bmsection*{Financial disclosure}

None reported.

\bmsection*{Conflict of interest}

The authors declare no potential conflict of interests.

\bmsection*{Data availability statement}

% Cut 2026-10-15 (D-20261015-1, needs Tak's veto): the IST-1 and IST-3 DataShare DOIs and 'the survey cited in Section 5' -> dropped with the survey (insurance, Louis: the registered survey is named as not reported; protocol and outputs stay with the code); the Study 2 implementation sentence (Rcpp, PSOCK, L'Ecuyer-CMRG, seed 42 + 1,000k, not bit-reproducible) moved here from Sec. 3.2.1
R code for computing the Wasserstein-1 distance and reproducing the simulation studies and the application is available at [repository URL]. The GUSTO-I individual patient data analysed in Section~\ref{sec:application} are publicly available as the dataset \texttt{gusto} in the R package predtools (version 0.0.3; \url{https://github.com/resplab/predtools}); the trial is described in the original publication;\cite{gusto1993} the analysis scripts load the data frame with \texttt{data(gusto, package = "predtools")} or from the identical CSV copy distributed with the code ($N = \numval{gustoN}$ in both). Study~2 used a PSOCK cluster with L'Ecuyer--CMRG random-number streams seeded per cell (seed $42 + 1{,}000k$ for cell $k = 1, \ldots, 21$); replicates are assigned to workers dynamically and are not bit-reproducible, so the summaries are reproducible to Monte Carlo error. A registered survey of region pairs in GUSTO-I, IST-1 and IST-3 is reported in Section~
ef{sec:discussion} only as the range of its rank correlations, its tables not being reproduced; its protocol and outputs are included with the R code, as are the registered protocols of the Study~2 extension and of the Set 4 control check and their outputs.

\bibliography{per_em_W1_wiley}

\bmsection*{Supporting information}

No supporting information accompanies this article; the R code, the registered protocols and all outputs are available as described in the Data availability statement.

%==============================================================================
\appendix
%==============================================================================

% Cut 2026-10-15 (D-20261015-1, needs Tak's veto): Appendix A.1 (non-negativity and identity of indiscernibles: textbook, Villani cited in Sec. 2.2) deleted; the Jensen derivation of the upward bias moved verbatim into the first appendix subsection, which now carries label app:asymptotic (the pointer to the old A.1 label in Sec. 2.3 is repointed); the asymptotics subsection shortened: the rate statement ('converges at rate O(n^{-1/2}) if and only if the integral of sqrt(F(1-F)) is finite') dropped, the grid-recorded sentences condensed -> deleted (no replacement wording)
\bmsection{Proofs and Mathematical Details\label{app:proofs}}

\bmsubsection{Upward Bias and Asymptotic Properties\label{app:asymptotic}}

Upward bias of $\widehat{W}_1$. Since $E\,\hat{F}_r(x) = F_r(x)$, Jensen's inequality gives $E\,|\hat{F}_1(x) - \hat{F}_2(x)| \geq |F_1(x) - F_2(x)|$ at every $x$, and integrating (the integrand is non-negative, so no further condition is needed) $E\,\widehat{W}_1 \geq W_1$ for all $n_1, n_2$.

Under standard regularity conditions, the empirical Wasserstein-1 distance $\widehat{W}_1 = W_1(\hat{F}_1, \hat{F}_2)$ is a consistent estimator of $W_1(F_1, F_2)$; when $\int \sqrt{F(1-F)}\,dx < \infty$, $\sqrt{n}\,W_1(\hat{F}_n, F)$ converges in distribution to $\int |B(F(x))|\,dx$ with $B$ a Brownian bridge,\cite{delbarrio1999} whose quantiles depend on the unknown $F$, so no universal critical values are available. For two samples with $n_1/(n_1+n_2) \to \lambda \in (0,1)$, when $\int \sqrt{F_r(1-F_r)}\,dx < \infty$ ($r = 1, 2$) and $\{x : F_1(x) = F_2(x) \in (0,1)\}$ has Lebesgue measure zero, the directional derivative of $W_1$ at $(F_1, F_2)$ is linear, which is the condition under which the bootstrap statement below holds;\cite{fang2019} this condition excludes $F_1 = F_2$. At the boundary $F_1 = F_2$ the derivative is non-linear and the limit is not Gaussian; separately, a percentile interval cannot contain zero, which is the degenerate coverage observed under the null (S1; Section~\ref{sec:sim_results}). For an effect modifier recorded on a grid, the CDFs are step functions and the condition fails whenever they agree at any grid point inside the support; the finite-space results of Sommerfeld and Munk then apply,\cite{sommerfeld2018} with a Gaussian limit when the dual solution is unique. The condition concerns the population CDFs and cannot be checked from the samples, so for a distance unresolved against its null floor, such as R2 on systolic blood pressure (Section~\ref{sec:app_operability}), $F_1 = F_2$ is not excluded and the interval is read at its upper limit only. Under the same conditions the percentile bootstrap confidence interval achieves asymptotically correct coverage;\cite{fang2019} where the two CDFs coincide on an interval, including $F_1 = F_2$, the naive bootstrap is not consistent.\cite{dumbgen1993,fang2019}

\bmsubsection{Comparison with the Kolmogorov--Smirnov Bound; Attainment on the Outcome Scale\label{app:ks_bound}}

% Moved from Section 2.2 on 2026-10-12 (D-20261012-3, EXT Minor 5): the technical comparison of the two bounds.
Where both constants are available, neither bound dominates in general. On a support $[a, b]$, $W_1 \leq (b - a)\,D_{\text{KS}}$ and $\mathrm{TV}(\tau) \leq L_{\text{clinical}}\,(b - a)$, and $L_{\text{clinical}} \cdot W_1 \leq \mathrm{TV}(\tau) \cdot D_{\text{KS}}$ holds exactly when $W_1 / D_{\text{KS}} \leq \mathrm{TV}(\tau) / L_{\text{clinical}}$. For a CATE whose slope has constant magnitude $L_{\text{clinical}}$, such as a linear CATE, $\mathrm{TV}(\tau) = L_{\text{clinical}}(b - a)$ and equation~(\ref{eq:heterogeneity_bound}) is never the looser of the two; the Kolmogorov--Smirnov bound gains as the CATE's slope falls below its maximum over much of the support, most sharply when the CATE changes steeply over a narrow range---the threshold case in which equation~(\ref{eq:heterogeneity_bound}) is itself conservative (Section~\ref{sec:discussion}). Every valid bound with the constant $L_{\text{clinical}}$, including $L_{\text{clinical}} \cdot W_2$ and, on $[a, b]$, $L_{\text{clinical}} (b - a) D_{\text{KS}}$, is at least $L_{\text{clinical}} W_1$, and the $L^*$ of Section~\ref{sec:inference} is the smallest slope at which a CATE restricted only by its slope produces the margin.

% Moved from Section 2.2 on 2026-10-14 (D-20261014-1, needs Tak's veto), verbatim: the sharpness of the Kolmogorov--Smirnov bound within its own class, and the bounded-support / unbounded-support caveat with the Lee et al. comparison.
% Cut 2026-10-15 (D-20261015-1, needs Tak's veto): '(supports that include the records flagged there)' dropped (pure deletion; the records are described in Sec. 4.1) -> deleted (no replacement wording)
The Kolmogorov--Smirnov bound is likewise sharp within its own class, by a step CATE of height $\mathrm{TV}(\tau)$ at the point where $|F_1 - F_2|$ is largest. A bound on the magnitude of $\tau$ itself tightens equation~(\ref{eq:heterogeneity_bound}) further: on a support $[a, b]$ the attaining $\tau^*$ ranges over up to $L_{\text{clinical}} (b - a)$, which on the risk-difference scale is \numval{lipRangeAge} for age and \numval{lipRangeSbp} for systolic blood pressure over the recorded supports of Section~\ref{sec:application}, exceeding the highest 30-day mortality Lee et al.\ report for any age group (20.5\% above 75 years\cite{lee1995}), so that the attaining CATE would require effects far beyond the baseline risks; for distributions with unbounded support no CATE bounded on the outcome scale attains the bound at all. The 30-day mortality that Lee et al.\ report for GUSTO-I is 1.1\% below 45 years and 20.5\% above 75 years,\cite{lee1995} a rise of 19.4 percentage points over at least three decades and hence an average prognostic age gradient of at most about $6.5 \times 10^{-3}$ per year (Section~\ref{sec:inference}).

% Cut 2026-10-15 (D-20261015-1, needs Tak's veto): Appendix B (R code for the W1 estimator and its bootstrap interval, 25-line listing, lst:nabcd) deleted; the code is in the repository named in the Data availability statement -> deleted (no replacement wording)

\end{document}
